\SourcePage{323}
\SGAchapter{XVI}{幺幂秩为零的群}{Groupes de rang unipotent nul}
\begin{center}M. Raynaud\footnote{参见第 XV 讲第 1 页的注。}\end{center}
\footnotetext[0]{原文版本注（0）：2008 年 12 月 1 日 xy 版。}
\SGAsection{1}{一个浸入判据}
\textbf{1.1. 不是\Term{immersion}{浸入}的\Term{group-prescheme}{群预概形}\Term{monomorphism}{单态射}的例子。}我们将构造一个\Term{prescheme}{预概形} $S$、两个 $S$-\Term{group-prescheme}{群预概形} $G$ 和 $H$，以及一个不是\Term{immersion}{浸入}的 $S$-群\Term{monomorphism}{单态射} $u:G\to H$。群 $G$ 和 $H$ 在 $S$ 上将是\Term{finite-presentation}{有限表示}且\Term{flat}{平坦}的，$G$ 甚至在 $S$ 上\Term{smooth}{光滑}（参见第 VIII 讲 7.1 之前的一段及注（$*$），另见第 XVII 讲附录 III，4）。

先取 $S$ 为一个混合\Term{characteristic}{特征}、\Term{residue-characteristic}{剩余特征}为 $2$ 的\Term{discrete-valuation-ring}{离散赋值环} $A$ 的\Term{spectrum}{谱}。记 $t$ 为 $S$ 的\Term{generic-point}{泛点}，$s$ 为\Term{closed-point}{闭点}。

 a）令 $G$ 为从\Term{constant-group}{常值群} $G'=(\mathbb Z/2\mathbb Z)_S$ 中删去\Term{closed-fiber}{闭纤维} $(\mathbb Z/2\mathbb Z)_s$ 中不同于单位元的点所得的\Term{open-subgroup}{开子群}；令 $H$ 为二次\Term{root-of-unity}{单位根}组成的\Term{multiplicative-type}{乘法型}群 $(\mu_2)_S$（第 I 讲 4.4.4）。由 $S$ 的选取，$H_t$ \Term{isomorphism}{同构}于 $(\mathbb Z/2\mathbb Z)_t$，而 $H_s$ 是\Term{radicial}{根性}群。$H$ 的\Term{section}{截面} $(-1)$ 显然定义一个\Term{morphism}{态射} $G'\to H$，从而得到\Term{morphism}{态射} $u:G\to H$；它是\Term{monomorphism}{单态射}，却不是\Term{immersion}{浸入}（否则它必为\Term{isomorphism}{同构}，参见第 VIII 讲 §7）。这里第 VIII 讲 7.9 不适用，因为 ${}_2G=G$ 在 $S$ 上不\Term{finite}{有限}。

 b）记 $K$ 为从\Term{unit-group}{单位群}出发，将\Term{closed-fiber}{闭纤维}的唯一点“复制为两个”所得的\Term{etale}{平展}而不\Term{separated}{分离}的 $S$-\Term{group-prescheme}{群预概形}。取 $H$ 为积群 $(\mu_2)_S\times_S K$。记 $a$（相应地 $b$）为 $\mu_2$（相应地 $K$）在 $S$ 上唯一的不同于单位\Term{section}{截面}的\Term{section}{截面}，$h=(a,b)$ 为 $H$ 中相应的\Term{section}{截面}。给定 $h$ 就定义了一个 $S$-\Term{group-morphism}{群态射}
\[
u:G=(\mathbb Z/2\mathbb Z)_S\longrightarrow H.
\]
\SourcePage{324}
显然 $u$ 是\Term{monomorphism}{单态射}，而不是\Term{immersion}{浸入}。这里第 VIII 讲 7.9 不适用，因为 ${}_2H=H$ 在 $S$ 上不\Term{separated}{分离}。

 c）下面的例子更有意思，其中 $G$ 是具有\Term{connected}{连通}\Term{fiber}{纤维}的\Term{smooth}{光滑}群，甚至 $G=(\mathbb G_a)_S$（参见第 VIII 讲 7.10）。

取 $S$ 为\Term{equal-characteristic}{等特征} $2$ 的\Term{discrete-valuation-ring}{离散赋值环} $A$ 的\Term{spectrum}{谱}，记 $\pi$ 为 $A$ 的一个\Term{uniformizer}{一致化元}，$\mathfrak m$ 为 $A$ 的\Term{maximal-ideal}{极大理想}。考虑\Term{additive-group}{加法群} $(\mathbb G_a)_S$ 和环为 $A[X,Y]$ 的积群 $(\mathbb G_a\times\mathbb G_a)_S$，令 $G_1$ 为方程
\[X^2+X-\pi Y=0\tag{1}\]
定义的\Term{closed-subgroup}{闭子群}。群 $G_1$ 在 $S$ 上\Term{smooth}{光滑}，其\Term{generic-fiber}{泛纤维}\Term{isomorphism}{同构}于 $(\mathbb G_a)_t$，\Term{closed-fiber}{闭纤维}\Term{isomorphism}{同构}于 $(\mathbb G_a)_s$ 与 $(\mathbb Z/2\mathbb Z)_s$ 的积。因子 $(\mathbb Z/2\mathbb Z)_s$ 可取为 $(G_1)_s$ 中由方程
\[Y=X^2+X=0\tag{2}\]
定义的\Term{subgroup}{子群} $N_s$。

设 $N'$ 和 $N''$ 是 $G_1$ 的两个\Term{group-subscheme}{群子概形}，均\Term{isomorphism}{同构}于 $(\mathbb Z/2\mathbb Z)_S$，它们的\Term{closed-fiber}{闭纤维}均为 $N_s$，而\Term{generic-fiber}{泛纤维}互不相同。因此，群 $N'$ 和 $N''$ 分别由 $G_1$ 在 $S$ 上的两个\Term{section}{截面}给定，其坐标 $(a',b')$ 和 $(a'',b'')$ 满足
\[
\begin{cases}
 a'^2+a'-\pi b'=a''^2+a''-\pi b''=0,\\
 a'\equiv a''\equiv1\pmod{\mathfrak m},\qquad a'\ne a'',\\
 b'\equiv b''\equiv0\pmod{\mathfrak m}.
\end{cases}\tag{3}
\]
（例如可取 $(1,0)$ 和 $(1+\pi^2,\pi+\pi^3)$。）\Term{quotient-group}{商群} $G'=G_1/N'$ 和 $G''=G_1/N''$ 是\Term{representable}{可表示}的（第 V 讲 4.1）。记 $u$ 为\Term{canonical}{典范}\Term{morphism}{态射}
\[u:G_1\longrightarrow G'\times_S G'',\]
并令 $H$ 为 $G'\times_S G''$ 中等于 $u_t(G_1)_t$ 在 $G'\times_S G''$ 中的\Term{schematic-closure}{概形闭包}的\Term{group-subscheme}{群子概形}，于是 $u$ 通过 $H$ \Term{resolution}{分解}。$u$ 的核为 $K=N'\cap N''$；其\Term{generic-fiber}{泛纤维}为\Term{unit-group}{单位群}，\Term{closed-fiber}{闭纤维}为 $N_s$；特别地，$K$ 在 $S$ 上不\Term{flat}{平坦}。因此，在\Term{generic-fiber}{泛纤维}上，$u_t$ 是从 $(G_1)_t$ 到 $H_t$ 的\Term{isomorphism}{同构}。另一方面，我们断言 $H_s$ 不\Term{smooth}{光滑}。事实上，若 $H_s$ \Term{smooth}{光滑}，由于 $u_s$ 的核\Term{finite}{有限}，且 $(G_1)_s$ 与 $H_s$ 具有相同\Term{dimension}{维数}（即 $1$），$u_s$ 就是\Term{flat}{平坦}\Term{morphism}{态射}；又因 $G_1$ 和 $H$ 都\Term{flat}{平坦}，$u$ 便是\Term{flat}{平坦}\Term{morphism}{态射}，从而 $\operatorname{Ker}u$ 在 $S$ 上\Term{flat}{平坦}，矛盾。显然，$u$ 在 $G_1$ 的\Term{identity-component}{单位分支} $G$ 上的\Term{restriction}{限制}是\Term{monomorphism}{单态射}（$K\cap G$ 的\Term{fiber}{纤维}都是\Term{unit-group}{单位群}，故 $K\cap G$ 是\Term{unit-group}{单位群}，第 VI\textsubscript{B} 讲 2.9），但不是\Term{immersion}{浸入}（否则便是\Term{open-immersion}{开浸入}，从而 $H_s$ \Term{smooth}{光滑}）。这里第 VIII 讲 7.9 不适用，因为 ${}_2G=G$ 在 $S$ 上不\Term{finite}{有限}。

对于喜欢方程的读者，上例中的 $H$ 可以取为 $(\mathbb G_a\times_S\mathbb G_a)_S$ 的\Term{closed-subgroup}{闭子群}，其环为 $A[V,W]/(F)$，其中
\[
F=(a''b'-b''a')(a''^2V-a'^2W)-(a''V-a'W)^2
\]
\SourcePage{325}
（$a',a'',b',b''$ 满足（3））。$G$ 取为环为 $A[T]$ 的群 $(\mathbb G_a)_S$，\Term{morphism}{态射} $u:G\to H$ 取为下列映射定义的 $S$-\Term{morphism}{态射}：
\[V\longmapsto a'(a'T+\pi T^2),\qquad W\longmapsto a''(a''T+\pi T^2).\]
\begin{SGAresult}{注记}{1.2}
上述构造受到 Koizumi--Shimura 方法的启发\footnote{编者注：S. Koizumi 与 G. Shimura，\textit{Specialization of abelian varieties}，\textit{Sci. Papers Coll. Gen. Ed. Univ. Tokyo} \textbf{9}（1959），187--211。}。当 $S$ 为混合\Term{characteristic}{特征}环的\Term{spectrum}{谱}时，它不适用，因为这时 $G$ 中\Term{finite-order-point}{有限阶点}“过于\Term{rigidity}{刚性}”。
\end{SGAresult}
\textbf{1.3. \Term{immersion}{浸入}判据的陈述。}
\begin{SGAresult}{\Term{theorem}{定理}}{1.3}
设 $S$ 为\Term{locally-noetherian}{局部诺特}\Term{prescheme}{预概形}，$G$ 为一个在 $S$ 上\Term{smooth}{光滑}、\Term{finite-type}{有限型}、并且对 \Term{fpqc-topology}{fpqc 拓扑}\Term{local-scheme}{局部}具有\Term{cartan-subgroup}{Cartan 子群}的 $S$-\Term{group-prescheme}{群预概形}（第 XV 讲 6.1 和 7.3（i）），$H$ 为\Term{locally-finite-type}{局部有限型}的 $S$-\Term{group-prescheme}{群预概形}，$u:G\to H$ 为一个 $S$-群\Term{monomorphism}{单态射}。则：

 a）若 $G$ 的\Term{fiber}{纤维}\Term{connected}{连通}，则 $u$ 为\Term{immersion}{浸入}的（必要且）充分条件是：对于每个作为\Term{residue-field}{剩余域}\Term{algebraically-closed}{代数闭}的\Term{complete}{完备}\Term{discrete-valuation-ring}{离散赋值环}之\Term{spectrum}{谱}的 $S$-\Term{scheme}{概形} $S'$，以及 $G_{S'}$ 的每个\Term{cartan-subgroup}{Cartan 子群} $C$，$u_{S'}$ 在 $C$ 上的\Term{restriction}{限制}都是\Term{immersion}{浸入}。

 b）$u$ 为\Term{immersion}{浸入}（相应地，\Term{closed-immersion}{闭浸入}）的（必要且）充分条件是：对于上述每个 $S'$，以及 $G_{S'}$ 的\Term{identity-component}{单位分支} $(G_{S'})^0$ 的每个\Term{cartan-subgroup}{Cartan 子群} $C$，$u_{S'}$ 在 $C$ 于 $G_{S'}$ 中的\Term{normalizer}{正规化子} $N$（第 XI 讲 6.11）上的\Term{restriction}{限制}都是\Term{immersion}{浸入}（相应地，\Term{closed-immersion}{闭浸入}）。
\end{SGAresult}
在证明 1.3 之前，先陈述几个应用。
\begin{SGAresult}{推论}{1.4}
设 $S$ 为\Term{prescheme}{预概形}，$G$ 为在 $S$ 上\Term{smooth}{光滑}、\Term{finite-presentation}{有限表示}、\Term{unipotent-rank}{幺幂秩}（第 XV 讲 6.1 ter）为零，并且对 \Term{fpqc-topology}{fpqc 拓扑}\Term{local-scheme}{局部}具有\Term{maximal-torus}{极大环面}的 $S$-\Term{group-prescheme}{群预概形}，$H$ 为 $S$-\Term{group-prescheme}{群预概形}，$u:G\to H$ 为 $S$-群\Term{monomorphism}{单态射}。再假设或者 $H$ 在 $S$ 上\Term{finite-presentation}{有限表示}，或者 $S$ \Term{locally-noetherian}{局部诺特}且 $H$ \Term{locally-finite-type}{局部有限型}。则：

 a）若 $G$ 的\Term{fiber}{纤维}\Term{connected}{连通}，则 $u$ 是\Term{immersion}{浸入}。

 b）若 $H$ 在 $S$ 上\Term{separated}{分离}，并且对于 $S$ 上的每个 $S'$ 和 $G_{S'}$ 的每个\Term{maximal-torus}{极大环面} $T$，\Term{weyl-group}{Weyl 群}
\[W=\operatorname{Norm}_{G_{S'}}(T)/\operatorname{Centr}_{G_{S'}^0}(T)\]
是\Term{representable}{可表示}的（当 $G$ 的\Term{fiber}{纤维}\Term{connected}{连通}时，此条件总成立，第 XV 讲 7.1（iv）），且在 $S$ 上\Term{finite}{有限}，则 $u$ 是\Term{closed-immersion}{闭浸入}。
\end{SGAresult}
\begin{SGAresult}{推论}{1.5}
设 $S$ 为\Term{prescheme}{预概形}，$G$ 为在 $S$ 上\Term{smooth}{光滑}、\Term{finite-presentation}{有限表示}且\Term{fiber}{纤维}\Term{connected}{连通}的 $S$-\Term{group-prescheme}{群预概形}，$H$ 为 $S$-\Term{group-prescheme}{群预概形}，$u:G\to H$ 为 $S$-群\Term{monomorphism}{单态射}。假设或者 $H$ 在 $S$ 上\Term{finite-presentation}{有限表示}，或者 $S$ \Term{locally-noetherian}{局部诺特}且 $H$ \Term{locally-finite-type}{局部有限型}。则：
\SourcePage{326}
 a）若 $G$ 为\Term{reductive-group}{约化群}，即在 $S$ 上\Term{affine}{仿射}且\Term{fiber}{纤维}\Term{reductive}{约化}（参见第 XIX 讲），则 $u$ 为\Term{immersion}{浸入}；若 $H$ \Term{separated}{分离}，则为\Term{closed-immersion}{闭浸入}。

 b）若 $G$ 的\Term{fiber}{纤维}\Term{affine}{仿射}、\Term{solvable}{可解}、\Term{connected}{连通}且\Term{unipotent-rank}{幺幂秩}为零，则 $u$ 是\Term{immersion}{浸入}；若 $H$ \Term{separated}{分离}，则为\Term{closed-immersion}{闭浸入}。
\end{SGAresult}
\textbf{由 1.3 \Term{pushout}{推出} 1.4 的证明。}

\emph{归约到 $S$ \Term{noetherian}{诺特}的情形。}若 $S$ \Term{locally-noetherian}{局部诺特}，由于待证性质在 $S$ 上是\Term{local-scheme}{局部}的，归约是直接的。在第二种情形中，$G$ 和 $H$ 在 $S$ 上\Term{finite-presentation}{有限表示}。缩小 $S$ 后，可以假设 $S$ \Term{affine}{仿射}。由第 VI\textsubscript{B} 讲 §10，存在\Term{noetherian}{诺特}\Term{scheme}{概形} $S_0$、在 $S_0$ 上\Term{smooth}{光滑}（在 a）中还具有\Term{connected}{连通}\Term{fiber}{纤维}）且\Term{finite-type}{有限型}的 $S_0$-\Term{group-prescheme}{群预概形} $G_0$、\Term{finite-type}{有限型}的 $S_0$-\Term{group-prescheme}{群预概形} $H_0$（在 b）中还\Term{separated}{分离}），以及一个 $S_0$-\Term{group-morphism}{群态射}
\[u_0:G_0\longrightarrow H_0,\]
使得 $G,H,u$ 由 $G_0,H_0,u_0$ 经过底\Term{group-extension}{扩张} $S\to S_0$ 得到。$u$ 为\Term{monomorphism}{单态射}\Term{equivalence}{等价}于 $\operatorname{Ker}u$ 为\Term{unit-group}{单位群}；因而可假设 $u_0$ 是\Term{monomorphism}{单态射}。由于 $G$ 的\Term{unipotent-rank}{幺幂秩} $\rho_u$ 为零，且对 \Term{fpqc-topology}{fpqc 拓扑}\Term{local-scheme}{局部}具有\Term{maximal-torus}{极大环面}，$G$ 的\Term{fiber}{纤维}的\Term{abelian-rank}{阿贝尔秩} $\rho_{\mathrm{ab}}$ \Term{locally-constant}{局部常值}（第 XV 讲 8.18）。而 $\rho_u$ 和 $\rho_{\mathrm{ab}}$ 是\Term{locally-constructible}{局部可构造}函数（第 XV 讲 6.3 bis）。一个标准论证（参见 EGA IV 8.3.4）说明，可选取 $S_0$ 和 $G_0$，使 $G_0$ 的\Term{fiber}{纤维}的\Term{unipotent-rank}{幺幂秩}为零、\Term{abelian-rank}{阿贝尔秩}\Term{locally-constant}{局部常值}。于是 $G_0$ 对 \Term{fpqc-topology}{fpqc 拓扑}\Term{local-scheme}{局部}具有\Term{maximal-torus}{极大环面}（第 XV 讲 8.18），而 $G_0$ 的\Term{maximal-torus}{极大环面}\Term{functor}{函子} $\mathscr{TM}_{G_0}$ 由一个在 $S_0$ 上\Term{finite-type}{有限型}的 $S_0$-\Term{scheme}{概形} $X_0$ 表示（第 XV 讲 8.15）。记 $T_0$ 为 $(G_0)_{X_0}$ 的“\Term{universal}{普遍}”\Term{maximal-torus}{极大环面}，$X$ 为 $S$-\Term{prescheme}{预概形} $X_0\times_{S_0}S$，$T=T_0\times_{S_0}S$ 为 $G$ 的\Term{universal}{普遍}\Term{maximal-torus}{极大环面}。由假设，在 b）中，关于 $T$ 的\Term{weyl-group}{Weyl 群}\Term{representable}{可表示}且在 $X$ 上\Term{finite}{有限}。这两个性质与\Term{prescheme}{预概形}的\Term{projective-limit}{投射极限}相容（第 VI\textsubscript{B} 讲 10.1 iii）和 EGA IV 8.10.5）。因而可选取 $S_0$，使 $T_0$ 的\Term{weyl-group}{Weyl 群}在 $X_0$ 上\Term{finite}{有限}。显然，这时为证明 1.4，可用 $S_0,G_0,u_0,H_0$ 替换 $S,G,u,H$，从而假设 $S$ \Term{noetherian}{诺特}。

使用 1.3 给出的\Term{valuative-criterion}{赋值判据}，便归约到 $S$ 是\Term{discrete-valuation-ring}{离散赋值环}之\Term{spectrum}{谱}且 $G^0$ 具有一个\Term{cartan-subgroup}{Cartan 子群} $C$ 的情形。

\emph{1.4 a）的证明。}必须证明 $u$ 在 $C$ 上的\Term{restriction}{限制}是\Term{immersion}{浸入}（1.3 a）），于是归约到 $G=C$ 的\Term{fiber}{纤维}\Term{connected}{连通}且\Term{nilpotent}{幂零}的情形。标准归约（参见第 VIII 讲 7.1 的证明）进一步将我们归约到 $H$ 在 $S$ 上\Term{flat}{平坦}、$u$ 在\Term{generic-fiber}{泛纤维}上为\Term{isomorphism}{同构}、在\Term{closed-fiber}{闭纤维}的底\Term{space}{空间}上为\Term{isomorphism}{同构}的情形。群 $H$ 此时具有\Term{connected}{连通}\Term{fiber}{纤维}，因而在 $S$ 上\Term{separated}{分离}（第 VI\textsubscript{B} 讲 5.2）。暂且承认以下\Term{lemma}{引理}。
\begin{SGAresult}{\Term{lemma}{引理}}{1.6}
设 $S$ 为\Term{prescheme}{预概形}，$G$ 为在 $S$ 上\Term{smooth}{光滑}、\Term{fiber}{纤维}\Term{connected}{连通}、\Term{nilpotent}{幂零}且\Term{unipotent-rank}{幺幂秩}为零的 $S$-\Term{group-prescheme}{群预概形}。则：

 i）$G$ 是\Term{commutative-group}{交换群}。
\SourcePage{327}
 ii）对每个整数 $n>0$，${}_nG$ 是在 $S$ 上\Term{flat}{平坦}且\Term{quasi-finite}{拟有限}的\Term{group-prescheme}{群预概形}；它在 $S$ 上\Term{finite}{有限}，当且仅当 $G$ 的\Term{abelian-rank}{阿贝尔秩}或\Term{reductive-rank}{约化秩}在 $S$ 上\Term{locally-constant}{局部常值}。

 iii）对每个在 $S$ 上\Term{invertible}{可逆}的整数 $q>0$，\Term{subgroup}{子群}族 ${}_{q^n}G$ 相对于 $S$ 在 $G$ 中\Term{universally-schematically-dense}{普遍概形稠密}（EGA IV 11.10.8）。
\end{SGAresult}
\Term{lemma}{引理} 1.6 适用于群 $G$；又因 $G$ \Term{local-scheme}{局部}具有\Term{maximal-torus}{极大环面}，其\Term{fiber}{纤维}的\Term{reductive-rank}{约化秩}\Term{locally-constant}{局部常值}，故由 1.6 ii），对每个整数 $n>0$，${}_nG$ 在 $S$ 上\Term{finite}{有限}。于是，由第 VIII 讲 7.9，$u$ 是\Term{immersion}{浸入}。

\textbf{1.6 的证明。}先考察 $S$ 为域之\Term{spectrum}{谱}的情形。
\begin{SGAresult}{\Term{lemma}{引理}}{1.7}
设 $k$ 为\Term{characteristic}{特征}为 $p$ 的域，$G$ 为\Term{smooth}{光滑}、\Term{connected}{连通}、\Term{nilpotent}{幂零}且\Term{unipotent-rank}{幺幂秩}为零的 $k$-\Term{algebraic-group}{代数群}。则 $G$ 是\Term{commutative-group}{交换群}，为\Term{abelian-variety}{阿贝尔簇} $A$ 通过\Term{torus}{环面} $T$ 的\Term{extension}{扩张}。对每个整数 $n>0$，${}_nG$ 为\Term{finite}{有限}群，由一个秩为 $n^{\rho_r+2\rho_{\mathrm{ab}}}$ 的 $k$-\Term{algebra}{代数}定义，且当 $(n,p)=1$ 时\Term{etale}{平展}。若 $(q,p)=1$，则\Term{subgroup}{子群}族 ${}_{q^n}G$（$n\in\mathbb N$）在 $G$ 中\Term{schematically-dense}{概形稠密}。
\end{SGAresult}
\begin{proof}[1.7 的证明]
记 $Z$ 为 $G$ 的\Term{center}{中心}。群 $G/Z$ 是\Term{affine}{仿射}的（第 XII 讲 6.1），\Term{unipotent-rank}{幺幂秩}为零、\Term{smooth}{光滑}且\Term{connected}{连通}，所以是\Term{torus}{环面}（BIBLE，4，\Term{theorem}{定理} 4），于是 $G$ 为\Term{commutative-group}{交换群}（第 XII 讲 6.4）。设 $T$ 为 $G$ 的唯一\Term{maximal-torus}{极大环面}（第 XV 讲 3.4），由 Chevalley \Term{theorem}{定理}（Bourbaki 讨论班 1956/57，第 145 号），立即得到 $G$ 是\Term{abelian-variety}{阿贝尔簇} $A$ 通过 $T$ 的\Term{extension}{扩张}。对每个整数 $n>0$，$T$ 中的 $n$ 次幂映射是\Term{epimorphism}{满态射}，由此得到\Term{exact-sequence}{正合列}
\[0\longrightarrow {}_nT\longrightarrow {}_nG\longrightarrow {}_nA\longrightarrow0.\]
Weil 的一个经典\Term{theorem}{定理}（A. Weil，\textit{Variétés abéliennes et courbes algébriques}，§IX，\Term{theorem}{定理} 33，推论 1）告诉我们，${}_nA$ 是由一个秩为 $n^{2\rho_{\mathrm{ab}}}$ 的 $k$-\Term{algebra}{代数}定义的\Term{finite}{有限}群。由于 ${}_nT$ 是秩为 $n^{\rho_r}$ 的\Term{finite}{有限}群，便得到所述的 ${}_nG$ 的结构。

现令 $H$ 为包含所有 ${}_{q^n}G$ 的 $G$ 的最小\Term{closed-subscheme}{闭子概形}。由上文和第 XV 讲 4.6，$H$ 是\Term{smooth}{光滑}且\Term{connected}{连通}的\Term{subgroup}{子群}，故与 $G$ 同属上述类型。$H$ 中的 $n$ 次幂映射是\Term{epimorphism}{满态射}（因为 ${}_qH$ \Term{finite}{有限}），故有\Term{exact-sequence}{正合列}
\[0\longrightarrow {}_qH\longrightarrow {}_qG\longrightarrow {}_q(G/H)\longrightarrow0.\]
由此得 ${}_q(G/H)=0$，故 $G/H=0$。也就是说，\Term{subgroup}{子群} ${}_{q^n}G$ 在 $G$ 中\Term{schematically-dense}{概形稠密}。
\end{proof}
\textbf{续证 1.6。}i）为证明 $G$ 为\Term{commutative-group}{交换群}，用通常方法先归约到 $S$ \Term{noetherian}{诺特}，再归约到 $S$ 是一个\Term{closed-point}{闭点}为 $s$ 的\Term{local-ring}{局部环}之\Term{spectrum}{谱}。$G$ 的\Term{center}{中心}由 $G$ 的一个\Term{closed}{闭}\Term{group-subscheme}{群子概形} $Z$ 表示（第 XI 讲 6.11）。为证明 $Z=G$，只需证明模去 $S$ 的环的\Term{maximal-ideal}{极大理想}的每个幂后均有 $Z=G$（因为这时 $Z$ 为 $G$ 的\Term{open-subgroup}{开子群}，而 $G$ 的\Term{fiber}{纤维}\Term{connected}{连通}，故 $Z=G$）。这将我们归约到 $S$ 为\Term{local-artinian}{局部阿廷}的情形。
\SourcePage{328}
设 $q$ 是在 $S$ 上\Term{invertible}{可逆}的整数。对每个整数 $n$，${}_{q^n}G_s$ 是 $G_s$ 的一个\Term{etale}{平展}且\Term{central-subgroup}{中心}的\Term{multiplicative-type}{乘法型}\Term{subgroup}{子群}（1.7）。由于 $G$ \Term{smooth}{光滑}，${}_{q^n}(G_s)$ \Term{liftable}{可提升}为 $G$ 的一个\Term{etale}{平展}且\Term{central-subgroup}{中心}的 $S$-\Term{group-subscheme}{群子概形} $M_n$（第 XV 讲 1.2）。于是，\Term{flat}{平坦}\Term{subgroup}{子群}族 $M_n$（$n\in\mathbb N$）在 $G$ 中\Term{schematically-dense}{概形稠密}（1.7 和 EGA IV 11.10.9）。由于 $Z$ 在 $G$ 中\Term{closed}{闭}且包含所有 $M_n$，确有 $Z=G$。

1.6 ii）为证明 ${}_nG$ 在 $S$ 上\Term{flat}{平坦}且\Term{quasi-finite}{拟有限}，只需证明 $G$ 中的 $n$ 次幂映射是\Term{flat}{平坦}且\Term{quasi-finite}{拟有限}的\Term{morphism}{态射}。因 $G$ 在 $S$ 上\Term{flat}{平坦}，可归约到 $S$ 是域之\Term{spectrum}{谱}的情形（EGA IV 11.3.10），而在 1.7 的证明中已指出 $n$ 次幂映射是\Term{epimorphism}{满态射}。此外，${}_nG$ 在 $S$ 上\Term{separated}{分离}，因为 $G$ 在 $S$ 上\Term{separated}{分离}（第 VI\textsubscript{B} 讲 5.2）。于是，${}_nG$ 在 $S$ 上\Term{finite}{有限}，当且仅当 ${}_nG$ 的\Term{fiber}{纤维}是秩\Term{locally-constant}{局部常值}的\Term{finite}{有限}\Term{algebra}{代数}的\Term{spectrum}{谱}（归约到 $S$ 是 \Term{henselian-local-ring}{Hensel 局部环}之\Term{spectrum}{谱}后，这一点立即可见）。而由 1.7，此条件\Term{equivalence}{等价}于 $G$ 的\Term{fiber}{纤维}的\Term{abelian-rank}{阿贝尔秩}或\Term{reductive-rank}{约化秩}\Term{locally-constant}{局部常值}。

1.6 iii）为证明\Term{subgroup}{子群}族 ${}_{q^n}G$（$n\in\mathbb N$）在 $G$ 中\Term{universally-schematically-dense}{普遍概形稠密}，仍归约到 $S$ \Term{noetherian}{诺特}的情形。结合 1.7 和 1.6 b），只需应用 EGA IV 11.10.9。

\textbf{1.4 b）的证明。}我们已归约到 $S$ 是\Term{discrete-valuation-ring}{离散赋值环}之\Term{spectrum}{谱}且 $G^0$ 具有一个\Term{cartan-subgroup}{Cartan 子群} $C$ 的情形。令 $N=\operatorname{Norm}_G C=\operatorname{Norm}_G T$，其中 $T$ 为 $C$ 的唯一\Term{maximal-torus}{极大环面}（第 XII 讲 7.1 a）和 b））。由于 $H$ 在 $S$ 上\Term{separated}{分离}，由 1.6 ii）及第 VIII 讲 7.12，$u$ 在 $C$ 上的\Term{restriction}{限制}是\Term{closed-immersion}{闭浸入}。另一方面，要证明 $u$ 是\Term{closed-immersion}{闭浸入}，只需证明 $u|_N$ 是\Term{closed-immersion}{闭浸入}（1.3 b））。由假设，$W=N/C$ 由一个\Term{finite}{有限}\Term{group-scheme}{群概形}表示，将下一个\Term{lemma}{引理}应用于\Term{exact-sequence}{正合列}
\[0\longrightarrow C\longrightarrow N\longrightarrow W\longrightarrow1\]
即可得证（注意，\Term{proper}{固有}的\Term{immersion}{浸入}是\Term{closed-immersion}{闭浸入}）。
\begin{SGAresult}{\Term{lemma}{引理}}{1.8}
设 $S$ 为\Term{prescheme}{预概形}，$G$ 为在 $S$ 上\Term{finite-presentation}{有限表示}的 $S$-\Term{group-prescheme}{群预概形}，它是一个在 $S$ 上\Term{proper}{固有}且\Term{finite-presentation}{有限表示}的\Term{group-prescheme}{群预概形} $G''$ 通过一个在 $S$ 上\Term{finite-presentation}{有限表示}且\Term{flat}{平坦}的\Term{group-prescheme}{群预概形} $G'$ 的\Term{extension}{扩张}：
\[1\longrightarrow G'\longrightarrow G\longrightarrow G''\longrightarrow1.\]
另设 $H$ 是在 $S$ 上\Term{finite-presentation}{有限表示}（当 $S$ \Term{locally-noetherian}{局部诺特}时，也可为\Term{locally-finite-presentation}{局部有限表示}）的 $S$-\Term{group-prescheme}{群预概形}，$u:G\to H$ 是一个 $S$-\Term{group-morphism}{群态射}。则若 $u$ 在 $G'$ 上的\Term{restriction}{限制}\Term{proper}{固有}，$u$ 就\Term{proper}{固有}。
\end{SGAresult}
\begin{proof}[1.8 的证明]
照例归约到 $S$ \Term{noetherian}{诺特}的情形（第 VI\textsubscript{B} 讲 §10 和第 XV 讲 6.2），便可应用\Term{properness-valuative-criterion}{固有性赋值判据}（EGA II 7.3.8）。于是设 $S$ 是一个\Term{discrete-valuation-ring}{离散赋值环} $A$ 的\Term{spectrum}{谱}，其\Term{closed-point}{闭点}为 $s$，\Term{generic-point}{泛点}为 $t$。设 $\bar x\in G(t)$，$h\in H(S)$，且 $u_t(\bar x)=h(t)$。必须证明 $\bar x$ 来自 $G(S)$ 中唯一的元素 $x$。甚至只需在对\Term{discrete-valuation-ring}{离散赋值环} $A$ 作\Term{faithfully-flat}{忠实平坦}\Term{group-extension}{扩张}后证明 $x$ 的存在唯一性。记 $\bar y$ 为 $\bar x$ 在 $G''(t)$ 中的\Term{projection}{投影}。
\SourcePage{329}
由于 $G''$ 在 $S$ 上\Term{proper}{固有}，$\bar y$ 来自 $G''(S)$ 中唯一的元素 $y$。设 $X$ 为 $y$ 在 $G$ 中的\Term{inverse-image}{逆像}。\Term{prescheme}{预概形} $X$ 在 $S$ 上\Term{faithfully-flat}{忠实平坦}（因为 $G'$ 在 $S$ 上\Term{faithfully-flat}{忠实平坦}，\Term{morphism}{态射} $G\to G''$ 也\Term{faithfully-flat}{忠实平坦}，第 VI\textsubscript{B} 讲 §9）。对\Term{discrete-valuation-ring}{离散赋值环}作\Term{faithfully-flat}{忠实平坦}\Term{group-extension}{扩张}后，可假设 $X$ 在 $S$ 上具有一个\Term{section}{截面} $\ell$（EGA IV 14.5.8）。以 $\bar x\ell_t^{-1}$ 替换 $\bar x$，可假设 $\bar x\in G'(t)$。于是 $x$ 的存在唯一性由 $u|_{G'}$ \Term{proper}{固有}得到。
\end{proof}
\textbf{1.5 的证明。}a）我们将在第 XIX 讲中看到，若 $G$ 是\Term{reductive-group}{约化群}，则 $G$ 对\Term{faithfully-flat-topology}{忠实平坦拓扑}\Term{local-scheme}{局部}具有\Term{maximal-torus}{极大环面}，其\Term{weyl-group}{Weyl 群}\Term{finite}{有限}，\Term{unipotent-rank}{幺幂秩}为零。由 1.4 即得 a）。

 b）若 $G$ 的\Term{fiber}{纤维}\Term{affine}{仿射}且\Term{unipotent-rank}{幺幂秩}为零，则 $G$ 对 \Term{fpqc-topology}{fpqc 拓扑}\Term{local-scheme}{局部}具有\Term{maximal-torus}{极大环面}（第 XV 讲 8.18）。只需应用 1.4，并注意 $G$ 的\Term{fiber}{纤维}\Term{connected}{连通}、\Term{smooth}{光滑}、\Term{affine}{仿射}且\Term{solvable}{可解}，故其\Term{weyl-group}{Weyl 群}为\Term{unit-group}{单位群}（BIBLE，6，\Term{theorem}{定理} 1，推论 3）。

\textbf{1.3 a）的证明。}由于\Term{morphism}{态射} $u$ 已是\Term{monomorphism}{单态射}，为证明 $u$ 是\Term{immersion}{浸入}，只需证明 $u$ 在 $u(G)$ 的各点处\Term{proper}{固有}（EGA IV 15.7.1）；为此可使用\Term{local-properness-valuative-criterion}{局部固有性赋值判据}（EGA IV 15.7.5）。于是归约到 $S$ 是一个\Term{residue-field}{剩余域}\Term{algebraically-closed}{代数闭}的\Term{complete}{完备}\Term{discrete-valuation-ring}{离散赋值环}之\Term{spectrum}{谱}，其\Term{closed-point}{闭点}为 $s$，\Term{generic-point}{泛点}为 $t$。因 $G$ 在 $S$ 上\Term{flat}{平坦}，标准归约（参见第 VIII 讲 7.1 的证明）使我们可以假设 $H$ 在 $S$ 上\Term{flat}{平坦}，$u_t$ 是\Term{isomorphism}{同构}，$u_s$ 是底\Term{space}{空间}上的\Term{isomorphism}{同构}。此时，$u$ 是\Term{immersion}{浸入}\Term{equivalence}{等价}于 $u$ 是\Term{isomorphism}{同构}。

由假设，群 $G$ 对 \Term{fpqc-topology}{fpqc 拓扑}\Term{local-scheme}{局部}具有\Term{cartan-subgroup}{Cartan 子群}，因而可以谈论 $G$ 的\Term{regular-point}{正则点}组成的\Term{open-subset}{开集} $U$（第 XV 讲 7.3 i）和 ii））。
\begin{SGAresult}{\Term{lemma}{引理}}{1.9}
在上述假设下，要使 $u$ 为\Term{immersion}{浸入}，只需 $u|_U$ 为\Term{immersion}{浸入}。
\end{SGAresult}
\begin{proof}
事实上，$u|_U$ 为\Term{immersion}{浸入}，意即存在 $H$ 的\Term{open-subset}{开集} $V$ 及 $V$ 的闭集 $F$，使 $u|_U$ 通过 $F$ \Term{resolution}{分解}，并诱导 $U$ 到 $F$ 的\Term{isomorphism}{同构}。由于 $H_t$，从而 $V_t$，是\Term{irreducible}{不可约}的，必有 $V_t=F_t$。于是 $F_t$ 包含 $V_t$ 在 $V$ 中的\Term{schematic-closure}{概形闭包}；因 $H$ 在 $S$ 上\Term{flat}{平坦}，此\Term{closure}{闭包}等于 $V$。总之，$V=F$。因此 $u|_U$ 为\Term{open-immersion}{开浸入}，从而 $u$ 为\Term{open-immersion}{开浸入}（第 VI\textsubscript{B} 讲 2.6）。
\end{proof}
现在只需证明 $u|_U$ 是\Term{immersion}{浸入}；为此应用\Term{local-properness-valuative-criterion}{局部固有性赋值判据}。以一个在 $S$ 上\Term{faithfully-flat}{忠实平坦}的\Term{discrete-valuation-ring}{离散赋值环}之\Term{spectrum}{谱}替换 $S$ 后，必须证明：若 $h$ 是 $H$ 在 $S$ 上的一个\Term{section}{截面}，其像 $h(S)$ 包含在 $u(U)$ 中，则 $h$ 是 $U$ 在 $S$ 上某个\Term{section}{截面} $g$ 的像。只需证明 $h$ 包含在 $G$ 的某个\Term{cartan-subgroup}{Cartan 子群} $C$ 的像中。
\SourcePage{330}
事实上，由假设 $u|_C$ 是\Term{immersion}{浸入}，故 $h$ 是 $C$ 的某个\Term{section}{截面} $g$ 的像，而 $g$ 必为 $U$ 的\Term{section}{截面}。

设 $a\in U(s)$、$b\in U(t)$ 满足 $u_s(a)=h(s)$、$u_t(b)=h(t)$。由于 $u_t$ 是\Term{isomorphism}{同构}，$h(t)$ 是 $H_t$ 的\Term{regular-point}{正则点}，故包含在 $H_t$ 的唯一\Term{cartan-subgroup}{Cartan 子群} $D_t$ 中（第 XIII 讲 3.2）。记 $D$ 为 $D_t$ 在 $H$ 中的\Term{schematic-closure}{概形闭包}。
\begin{SGAresult}{\Term{lemma}{引理}}{1.10}
 i）$D_s$ 是\Term{nilpotent}{幂零}\Term{algebraic-group}{代数群}（第 VI\textsubscript{B} 讲 §8）。

 ii）$\dim D_s=\nu=G$ 的\Term{nilpotent-rank}{幂零秩} $=(H_s)_{\mathrm{red}}$ 的\Term{nilpotent-rank}{幂零秩}。

 iii）$h(s)\in D_s(\kappa(s))$。
\end{SGAresult}
\begin{proof}
 i）由于 $H$ \Term{separated}{分离}（第 VI\textsubscript{B} 讲 5.2），$D$ 也\Term{separated}{分离}。另一方面，$D$ 在 $S$ 上\Term{flat}{平坦}，其\Term{generic-fiber}{泛纤维}\Term{nilpotent}{幂零}。由第 VI\textsubscript{B} 讲 8.4，$D_s$ \Term{nilpotent}{幂零}。

 ii）由于 $(H_s)_{\mathrm{red}}\simeq G_s$，这两个群具有相同的\Term{nilpotent-rank}{幂零秩}。群 $G$ 对 \Term{fpqc-topology}{fpqc 拓扑}\Term{local-scheme}{局部}具有\Term{cartan-subgroup}{Cartan 子群}，故 $G_s$ 和 $G_t$ 的\Term{nilpotent-rank}{幂零秩}相同，记为 $\nu$。最后，群 $D$ 在 $S$ 上\Term{flat}{平坦}，故 $\dim D_s=\dim D_t=\nu$（第 VI\textsubscript{B} 讲 §4）。

 c）由于 $h_t$ 通过 $D_t$ \Term{resolution}{分解}，$h$ 显然通过 $D$ \Term{resolution}{分解}，从而得到 iii）。
\end{proof}
在此条件下，下一个\Term{lemma}{引理}证明 $(D_s)_{\mathrm{red}}$ 是 $(H_s)_{\mathrm{red}}\simeq G_s$ 的\Term{cartan-subgroup}{Cartan 子群}。
\begin{SGAresult}{\Term{lemma}{引理}}{1.11}
设 $G$ 为定义在\Term{algebraically-closed}{代数闭}域 $k$ 上的\Term{smooth}{光滑}\Term{connected}{连通}\Term{algebraic-group}{代数群}，$D$ 为一个\Term{smooth}{光滑}、\Term{nilpotent}{幂零}的\Term{algebraic-subgroup}{代数子群}，包含 $G(k)$ 的某个\Term{regular-element}{正则元素} $a$。则 $D^0$ 包含在 $G$ 的某个\Term{cartan-subgroup}{Cartan 子群}中。若再有 $\dim D$ 等于 $G$ 的\Term{nilpotent-rank}{幂零秩}，则 $D$ 是 $G$ 的\Term{cartan-subgroup}{Cartan 子群}（因而\Term{connected}{连通}）。
\end{SGAresult}
\begin{proof}
记 $Z$ 为 $G$ 的\Term{center}{中心}，$G'=G/Z$，它是\Term{affine}{仿射}的（第 XII 讲 6.1）；记 $a',D'$ 为 $a,D$ 在 $G'$ 中的像。由 $G$ 与 $G'$ 的\Term{cartan-subgroup}{Cartan 子群}之间的对应（第 XII 讲 6.6 e）），立即看出 $a'$ 是 $G'$ 的\Term{regular-element}{正则元素}，且只需对 $(D',G')$ 证明本\Term{lemma}{引理}；因此可以假设 $G$ \Term{affine}{仿射}。

现记 $s$ 为 $a$ 的\Term{semisimple}{半单}部分，它属于 $D(k)$（BIBLE，4，\Term{theorem}{定理} 3），且是\Term{regular-element}{正则元素}（BIBLE，7，\Term{theorem}{定理} 2，推论 1）。由 BIBLE，6，\Term{theorem}{定理} 2，$s$ 与 $D^0$ 中各元交换，故 $D^0$ 包含在 $s$ 的\Term{connected-centralizer}{连通中心化子}中，而后者是 $G$ 的一个\Term{cartan-subgroup}{Cartan 子群}（BIBLE，7，\Term{theorem}{定理} 2）。若现在 $\dim D$ 等于 $G$ 的\Term{nilpotent-rank}{幂零秩}，则 $D^0$ 就是 $G$ 的\Term{cartan-subgroup}{Cartan 子群}，等于 $\operatorname{Centr}_G(T)$，其中 $T$ 为 $D^0$ 的唯一\Term{maximal-torus}{极大环面}（第 XII 讲 6.6）。但 $D$ \Term{nilpotent}{幂零}，故与 $T$ 中各元交换（BIBLE，6，\Term{theorem}{定理} 2），从而 $D=D^0$。
\end{proof}
现在回到群 $G$。群 $C_t=u_t^{-1}(D_t)$ 是 $G_t$ 中包含 $b$ 的唯一\Term{cartan-subgroup}{Cartan 子群}。记 $C$ 为 $C_t$ 在 $G$ 中的\Term{schematic-closure}{概形闭包}。由\Term{schematic-closure}{概形闭包}的\Term{functoriality}{函子性}，$u|_C$ 通过 $D$ \Term{resolution}{分解}，故 $u_s(C_s)\subset D_s$。由于 $u_s$ 是 $G_s$ 到 $(H_s)_{\mathrm{red}}$ 的\Term{isomorphism}{同构}，且
\SourcePage{331}
$\dim C_s=\dim C_t=\nu=\dim D_s$（1.10），又有 $D_s$ \Term{connected}{连通}（1.11），$u_s$ 给出 $(C_s)_{\mathrm{red}}$ 到 $(D_s)_{\mathrm{red}}$ 的\Term{isomorphism}{同构}，故 $(C_s)_{\mathrm{red}}$ 是 $G_s$ 的\Term{cartan-subgroup}{Cartan 子群}。下一个\Term{lemma}{引理}说明，实际上 $C$ 是 $G$ 的\Term{cartan-subgroup}{Cartan 子群}。
\begin{SGAresult}{\Term{lemma}{引理}}{1.12}
设 $S$ 为\Term{discrete-valuation-ring}{离散赋值环}之\Term{spectrum}{谱}，$G$ 为\Term{smooth}{光滑}、\Term{finite-type}{有限型}的 $S$-\Term{group-prescheme}{群预概形}，$C$ 为在 $S$ 上\Term{flat}{平坦}的 $G$ 的\Term{group-subprescheme}{群子预概形}，使其\Term{generic-fiber}{泛纤维} $C_t$ 为 $G_t$ 的\Term{cartan-subgroup}{Cartan 子群}，而\Term{reduced}{既约}的几何\Term{closed-fiber}{闭纤维} $(C_{\bar s})_{\mathrm{red}}$ 为 $G_{\bar s}$ 的\Term{cartan-subgroup}{Cartan 子群}。则 $C$ 是 $G$ 的\Term{cartan-subgroup}{Cartan 子群}。
\end{SGAresult}
\begin{proof}
必须证明 $C$ 在 $S$ 上\Term{smooth}{光滑}，且只需在\Term{faithfully-flat}{忠实平坦}底\Term{group-extension}{扩张}后证明。关于 $C$ 的假设蕴涵 $G$ 的\Term{fiber}{纤维}的\Term{nilpotent-rank}{幂零秩}恒定，因此可谈论 $G$ 的\Term{regular-point}{正则点}\Term{open-subset}{开集} $U$（第 XV 讲 7.3）。由于 $(C_{\bar s})_{\mathrm{red}}$ 是 $G_{\bar s}$ 的\Term{cartan-subgroup}{Cartan 子群}，$C_{\bar s}$ 包含 $G_{\bar s}$ 的一个\Term{regular-point}{正则点} $a_{\bar s}$。但 $C$ 在 $S$ 上\Term{flat}{平坦}，故由 EGA IV 14.5.8，作\Term{faithfully-flat}{忠实平坦}底\Term{group-extension}{扩张}后，可假设 $a_{\bar s}$ 是 $G(s)$ 中的元素，并\Term{lifting}{提升}为 $C(S)$ 的\Term{section}{截面} $a$。$U$ 是包含 $a(s)$ 的\Term{open-subset}{开集}，故包含 $a$。令 $C'$ 为 $G$ 中包含 $a$ 的唯一\Term{cartan-subgroup}{Cartan 子群}（第 XV 讲 7.3）。必有 $C_t=C'_t$，因而 $C'$ 与 $C$ 均等于 $C_t$ 在 $G$ 中的\Term{schematic-closure}{概形闭包}的\Term{identity-component}{单位分支}（注意 $C'$ 和 $C$ 在 $S$ 上\Term{flat}{平坦}）。
\end{proof}
\textbf{结束 1.3 a）的证明。}由假设，$u$ 在 $G$ 的\Term{cartan-subgroup}{Cartan 子群} $C$ 上的\Term{restriction}{限制}是\Term{immersion}{浸入}。显然 $h(S)$ 包含在 $u(C)=D$ 中，故 $h$ 是 $C(S)$ 的一个\Term{section}{截面} $g$ 的像。证毕。

\textbf{1.3 b）的证明。}对\Term{immersion}{浸入}情形，仍使用\Term{local-properness-valuative-criterion}{局部固有性赋值判据}（EGA IV 15.5）；对\Term{closed-immersion}{闭浸入}情形，则使用\Term{properness-valuative-criterion}{固有性赋值判据}（EGA II 7.3.6）。于是归约到 $S$ 是\Term{residue-field}{剩余域}\Term{algebraically-closed}{代数闭}的\Term{complete}{完备}\Term{discrete-valuation-ring}{离散赋值环}之\Term{spectrum}{谱}。记 $s$ 为 $S$ 的\Term{closed-point}{闭点}，$t$ 为\Term{generic-point}{泛点}。以 $u_t(G_t)$ 在 $H$ 中的\Term{schematic-closure}{概形闭包}替换 $H$，可假设 $H$ 在 $S$ 上\Term{flat}{平坦}且 $u_t$ 为\Term{isomorphism}{同构}。

记 $G^0$ 为 $G$ 的\Term{connected-component}{连通分支}。由 1.2 a），$u|_{G^0}$ 是\Term{immersion}{浸入}；以 $H'$ 记 $G^0$ 的像所给出的 $H$ 的\Term{group-subprescheme}{群子预概形}。于是 $H'_s=(H_s)^0_{\mathrm{red}}$。

为验证\Term{valuative-criterion}{赋值判据}，必须证明：在\Term{immersion}{浸入}情形中，每个满足 $h(S)\subset u(G)$ 的 $H$ 在 $S$ 上的\Term{section}{截面} $h$（相应地，在\Term{closed-immersion}{闭浸入}情形中，每个 $H$ 在 $S$ 上的\Term{section}{截面} $h$），都是 $G$ 在 $S$ 上某个\Term{section}{截面} $g$ 的像。

设 $C$ 为 $G$ 的\Term{cartan-subgroup}{Cartan 子群}，$D=u(C)$ 为其像所给出的 $H'$ 的\Term{cartan-subgroup}{Cartan 子群}。$S$-群 $D'=\operatorname{int}(h)D$ 的\Term{fiber}{纤维}\Term{connected}{连通}，故其底\Term{space}{空间}包含在 $H'$ 的底\Term{space}{空间}中；另一方面，它在 $S$ 上\Term{smooth}{光滑}，故\Term{reduced}{既约}，所以 $D'$ 为 $H'$ 的一个\Term{smooth}{光滑}\Term{group-subprescheme}{群子预概形}。$D'$ 的\Term{fiber}{纤维}是 $H'$ 各\Term{fiber}{纤维}的\Term{cartan-subgroup}{Cartan 子群}，故 $D'$ 为 $H'$ 的\Term{cartan-subgroup}{Cartan 子群}，且 $C'=u^{-1}(D')$ 为 $G$ 的\Term{cartan-subgroup}{Cartan 子群}。而 $\operatorname{Transp}^{\mathrm{str}}_G(C,C')$ 在 $S$ 上\Term{smooth}{光滑}且\Term{surjectivity}{满射}（第 XII 讲 7.1 b））。由于 $S$ 为 Hensel 的且\Term{residue-field}{剩余域}\Term{algebraically-closed}{代数闭}，存在 $g\in G(S)$ 使 $\operatorname{int}(g)C=C'$。以 $u(g)^{-1}h$ 替换 $h$ 后，可以
\SourcePage{332}
假设 $h(t)$ \Term{normalization}{正规化} $D_t$（并且在\Term{immersion}{浸入}情形中，$h(s)$ 是 $C_s$ 在 $G_s$ 中的\Term{normalizer}{正规化子}的某个元素的像）。令 $N=\operatorname{Norm}_G(C)$。由假设，$u|_N$ 是\Term{immersion}{浸入}（相应地，\Term{closed-immersion}{闭浸入}），故 $h$ 确为 $N$ 在 $S$ 上某个\Term{section}{截面}的像，证明完毕。

\SGAsection{2}{商的一个可表示性定理}
“回顾”以下结果。
\begin{SGAresult}{\Term{theorem}{定理}}{2.1}
设 $S$ 为\Term{prescheme}{预概形}，$X$ 和 $Y$ 为两个 $S$-\Term{prescheme}{预概形}，$f:X\to Y$ 为 $S$-\Term{morphism}{态射}。假设下列两种情形之一成立：

 a）\Term{morphism}{态射} $f$ \Term{locally-finite-presentation}{局部有限表示}。

 b）\Term{prescheme}{预概形} $S$ \Term{locally-noetherian}{局部诺特}，而 $X$ 在 $S$ 上\Term{locally-finite-type}{局部有限型}。

则下列条件\Term{equivalence}{等价}：

 i）存在一个 $S$-\Term{prescheme}{预概形} $X'$ 及 $f$ 的\Term{resolution}{分解}
\[f:X\xrightarrow{f'}X'\xrightarrow{f''}Y,\]
其中 $f'$ 是\Term{faithfully-flat}{忠实平坦}且\Term{locally-finite-presentation}{局部有限表示}的 $S$-\Term{morphism}{态射}，而 $f''$ 是\Term{monomorphism}{单态射}。

 ii）（第一）\Term{projection}{投影}
\[p_1:X\times_Y X\longrightarrow X\]
是\Term{flat}{平坦}\Term{morphism}{态射}。

此外，若上述条件成立，则 $(X',f')$ 是 $X$ 关于 $f$ 所定义的\Term{equivalence-relation}{等价关系}的商（对 \Term{fpqc-topology}{fpqc 拓扑}而言），故 i）中的\Term{resolution}{分解} $f=f''\circ f'$ 在\Term{isomorphism}{同构}意义下唯一。
\end{SGAresult}
这个精细的\Term{theorem}{定理}的证明将在 EGA V 中给出；也可参阅 J.-P. Murre 在 Bourbaki 讨论班 1965 年 5 月第 294 号讲演的\Term{theorem}{定理} 2、推论 2，其中处理了 $Y$ \Term{locally-noetherian}{局部诺特}且 $X$ 在 $Y$ 上\Term{finite-type}{有限型}的情形。下面说明如何归约到此情形。

先作几条说明。

 a）i）$\Rightarrow$ ii）是显然的。事实上，第一\Term{projection}{投影}
\[p'_1:X\times_{X'}X\longrightarrow X\]
通过 $X\times_Y X$ \Term{resolution}{分解}为
\[p'_1:X\times_{X'}X\xrightarrow{u}X\times_Y X\xrightarrow{p_1}X.\]
由于 $f''$ 为\Term{monomorphism}{单态射}，\Term{morphism}{态射} $u$ 是\Term{isomorphism}{同构}；由于 $f'$ \Term{flat}{平坦}，$p'_1$ \Term{flat}{平坦}，从而 $p_1$ \Term{flat}{平坦}。
\SourcePage{333}
 b）2.1 的各项断言在 $Y$ 上是\Term{local-scheme}{局部}的（从而在 $S$ 上是\Term{local-scheme}{局部}的）；它们在 $X$ 上也是\Term{local-scheme}{局部}的，这容易由\Term{flat}{平坦}且\Term{locally-finite-presentation}{局部有限表示}的\Term{morphism}{态射}是开映射这一事实得到（EGA IV 11.3.1）。

 c）在 2.1 a）的假设下，由上文可归约到 $X$ 和 $Y$ \Term{affine}{仿射}且 $f$ \Term{finite-presentation}{有限表示}的情形。以 $Y$ 替换 $S$ 后，可假设 $X$ 和 $Y$ 在 $S$ 上\Term{finite-presentation}{有限表示}。然后借助 EGA IV 11.2.6 归约到 $S$ \Term{noetherian}{诺特}的情形。

 d）在 2.1 b）的假设下，可假设 $S,X,Y$ 都\Term{affine}{仿射}，$S$ \Term{noetherian}{诺特}，$X$ 在 $S$ 上\Term{finite-type}{有限型}。将 $Y$ 看作在 $S$ 上\Term{finite-type}{有限型}的\Term{affine}{仿射}\Term{scheme}{概形} $Y_i$ 的\Term{filtered}{滤过}\Term{projective-limit}{投射极限}。\Term{scheme}{概形} $X\times_{Y_i}X$ 构成 $X\times_S X$ 的\Term{closed-subscheme}{闭子概形}的\Term{filtered}{滤过}递减族，其\Term{projective-limit}{投射极限}为 $X\times_Y X$。由于 $X\times_S X$ \Term{noetherian}{诺特}，当 $i$ 充分大时，$X\times_{Y_i}X=X\times_Y X$，故若 $f$ 满足 2.1 ii）的假设，则 $f_i:X\xrightarrow{f}Y\to Y_i$ 也满足。由于 $f$ 在 $X$ 上定义的\Term{equivalence-relation}{等价关系}与 $f_i$ 定义的相同，显然只需对 $f_i$ 证明 ii）$\Rightarrow$ i），从而归约到 $Y$ 在 $S$ 上\Term{finite-type}{有限型}的情形。

\textbf{对\Term{group-prescheme}{群预概形}的应用。}
\begin{SGAresult}{\Term{theorem}{定理}}{2.2}
设 $S$ 为\Term{prescheme}{预概形}，$G$ 为在 $S$ 上\Term{locally-finite-presentation}{局部有限表示}并\Term{action}{作用}于一个 $S$-\Term{prescheme}{预概形} $X$ 的 $S$-\Term{group-prescheme}{群预概形}。设 $\xi\in X(S)$ 为 $X$ 在 $S$ 上的一个\Term{section}{截面}，使 $\xi$ 在 $G$ 中的\Term{stabilizer}{稳定子} $H$ 为 $G$ 的一个在 $S$ 上\Term{flat}{平坦}的\Term{group-subprescheme}{群子预概形}。则若 $X$ 在 $S$ 上\Term{locally-finite-type}{局部有限型}，或 $S$ \Term{locally-noetherian}{局部诺特}，商\Term{homogeneous-space}{齐性空间} $G/H$ 由一个在 $S$ 上\Term{locally-finite-presentation}{局部有限表示}的 $S$-\Term{prescheme}{预概形}表示，且 $S$-\Term{morphism}{态射}
\[f:G\longrightarrow X,\qquad g\longmapsto g\cdot\xi\]
\Term{resolution}{分解}为
\[
\begin{tikzcd}
 G\arrow[d,"p"']\arrow[dr,"f"] &\\
 G/H\arrow[r,"i"']&X,
\end{tikzcd}
\]
其中 $p$ 为\Term{canonical-projection}{典范投影}，它是\Term{faithfully-flat}{忠实平坦}且\Term{locally-finite-presentation}{局部有限表示}的\Term{morphism}{态射}，而 $i$ 为\Term{monomorphism}{单态射}。（为明确起见，我们假定 $G$ \Term{left-action}{左作用}于 $X$。）
\end{SGAresult}
\begin{proof}
\Term{morphism}{态射} $f$ 将 $G$ 赋予一个 $X$-\Term{prescheme}{预概形}结构。由 $\xi$ 的\Term{stabilizer}{稳定子}的定义，\Term{morphism}{态射}
\[G\times_S H\longrightarrow G\times_X G,\qquad(g,h)\longmapsto(g,gh)\]
是\Term{isomorphism}{同构}。由于 $H$ 在 $S$ 上\Term{flat}{平坦}，$G\times_S H$ 在 $G$ 上\Term{flat}{平坦}，故第一\Term{projection}{投影} $p_1:G\times_S G\to G$ 是\Term{flat}{平坦}\Term{morphism}{态射}。另一方面，若 $X$ 在 $S$ 上\Term{locally-finite-type}{局部有限型}，则 $f$ \Term{locally-finite-presentation}{局部有限表示}（EGA IV 1.4.3 v））；否则，假设 $S$ \Term{locally-noetherian}{局部诺特}。于是只需对\Term{morphism}{态射} $f$ 应用 2.1。
\SourcePage{334}
还需说明 $G/H$ 在 $S$ 上\Term{locally-finite-presentation}{局部有限表示}，这立即由第 V 讲 9.1 得到。
\end{proof}
\begin{SGAresult}{推论}{2.3}
设 $S$ 为\Term{prescheme}{预概形}，$G,H$ 为两个 $S$-\Term{group-prescheme}{群预概形}，$u:G\to H$ 为 $S$-\Term{homomorphism}{同态}。假设 $G$ 在 $S$ 上\Term{locally-finite-presentation}{局部有限表示}，且或者 $H$ 在 $S$ 上\Term{locally-finite-type}{局部有限型}，或者 $S$ \Term{locally-noetherian}{局部诺特}。则若 $K=\operatorname{Ker}(u)$ 在 $S$ 上\Term{flat}{平坦}，\Term{quotient-group}{商群} $G/K$ 由一个在 $S$ 上\Term{locally-finite-presentation}{局部有限表示}的 $S$-\Term{group-prescheme}{群预概形}表示，且 $u$ \Term{resolution}{分解}为
\[
\begin{tikzcd}
 G\arrow[rr,"u"]\arrow[dr,"p"']&&H\\
 &G/K\arrow[ur,"i"']&
\end{tikzcd}
\]
其中 $p$ 为\Term{canonical-projection}{典范投影}，$i$ 为\Term{monomorphism}{单态射}。
\end{SGAresult}
\begin{proof}
取 $X=H$，$\xi$ 为 $H$ 的单位\Term{section}{截面}，应用 2.2。
\end{proof}
\begin{SGAresult}{推论}{2.4}
设 $S$ 为\Term{prescheme}{预概形}，$G$ 为在 $S$ 上\Term{finite-presentation}{有限表示}的 $S$-\Term{group-prescheme}{群预概形}，$H$ 为在 $S$ 上\Term{smooth}{光滑}且\Term{fiber}{纤维}\Term{connected}{连通}的 $S$-\Term{group-prescheme}{群预概形}（因而由第 VI\textsubscript{B} 讲 5.5，它在 $S$ 上\Term{finite-presentation}{有限表示}），$i:H\to G$ 为 $S$-群的\Term{monomorphism}{单态射}。于是 $H$ 是 $G$ 的\Term{subgroup}{子群}。

假设 $N=\operatorname{Norm}_G(H)$（由第 XI 讲 6.11，它由 $G$ 的一个在 $S$ 上\Term{finite-presentation}{有限表示}的\Term{closed}{闭}\Term{group-subprescheme}{群子预概形}表示）在 $S$ 上\Term{flat}{平坦}。则 $G/N$ 由一个在 $S$ 上\Term{finite-presentation}{有限表示}且\Term{quasi-projective}{拟射影}的 $S$-\Term{prescheme}{预概形}表示。
\end{SGAresult}
\begin{proof}
除最后一项外，各项待证断言在 $S$ 上是\Term{local-scheme}{局部}的。为证明它们，可以假设 $S$ \Term{quasi-compact}{拟紧}，且 $H$ 在 $S$ 上的\Term{relative-dimension}{相对维数}恒等于 $r$。按第 XV 讲 §5 的方法进行。对每个整数 $n>0$，记 $G^{(n)}$（相应地 $H^{(n)}$）为 $G$（相应地 $H$）沿单位\Term{section}{截面}的第 $n$ 个\Term{normal-invariant}{法不变量}（invariant normal；EGA IV 16）。$\mathscr O_S$-模\Term{sheaf}{层} $H^{(n)}$ 是 $G^{(n)}$ 的商；由于 $H$ 在 $S$ 上\Term{smooth}{光滑}且\Term{dimension}{维数}为 $r$，$H^{(n)}$ \Term{canonical}{典范}地定义 $X_n=\operatorname{Grass}_{\varphi(n,r)}G^{(n)}$ 中的一个元素 $\xi_n$（EGA I，第 2 版，§9），其中 $\varphi(n,r)$ 为适当的整数。另一方面，$G$ 通过表示
\[G\longrightarrow\operatorname{Aut}_{S\text{-gr}}(G),\qquad g\longmapsto\operatorname{int}(g)\]
自然地\Term{action}{作用}于 $G^{(n)}$（从而也\Term{action}{作用}于 $X_n$）。由于 $S$ \Term{quasi-compact}{拟紧}，存在整数 $m$，使得当 $n\ge m$ 时有
\[N=\operatorname{Norm}_G(H)=\operatorname{Norm}_G H^{(n)}\qquad\text{（第 XI 讲 6.11 b））}.\]
即当 $n$ 充分大时，$N$ 是 $\xi_n$ 的\Term{stabilizer}{稳定子}。故 $G/N$ 的\Term{representability}{可表示性}由 2.2 得到。$G/N$ 在 $S$ 上\Term{finite-presentation}{有限表示}则由第 V 讲 9.1 得到。还需证明 $G/N$ 在 $S$ 上\Term{quasi-projective}{拟射影}；为此我们将给出 $G/N$ 上一个\Term{canonical}{典范}的 $S$-\Term{ample}{丰沛}\Term{invertible-sheaf}{可逆层}。考虑\Term{functor}{函子} $\mathscr F:(\operatorname{Sch}/S)^\circ\to\operatorname{Ens}$，对每个 $S$-\Term{prescheme}{预概形} $T$，规定
\[
\mathscr F(T)=\{\text{$G_T$ 的可表示、在 $T$ 上光滑且纤维连通的 $T$-子群 $H$}\}.
\]
\SourcePage{335}
这样的群 $H$ 在 $S$ 上\Term{finite-presentation}{有限表示}（第 VI\textsubscript{B} 讲 5.5），且 $H\to G_T$ 是\Term{quasi-affine}{拟仿射}\Term{morphism}{态射}（EGA IV 8.11.2）。由\Term{quasi-affine}{拟仿射}\Term{morphism}{态射}的\Term{effective-descent}{有效下降}（SGA 1，第 VIII 讲 7.9），可知 $\mathscr F$ 是 \Term{fpqc-topology}{fpqc 拓扑}的\Term{sheaf}{层}。由于 $G/N$ 是 $\mathscr F$ 的某个\Term{subsheaf}{子层}的\Term{associated-sheaf}{伴随层}，存在\Term{canonical}{典范}\Term{monomorphism}{单态射} $G/N\to\mathscr F$。所以存在 $G\times_S(G/N)$ 的\Term{subgroup}{子群} $H'$，它\Term{representable}{可表示}、在 $G/N$ 上\Term{smooth}{光滑}且\Term{fiber}{纤维}\Term{connected}{连通}，对于\Term{functor}{函子} $G/N$ 是“\Term{universal}{普遍}”的。我们断言\Term{invertible-sheaf}{可逆层} $\mathscr L=(\det(\operatorname{Lie}H'))^{-1}$ 是 $S$-\Term{ample}{丰沛}的。在这种表述下，断言在 $S$ 上是\Term{local-scheme}{局部}的，其证明与第 XV 讲 5.8 的证明类似。
\end{proof}
下面的推论已在第 XIV 讲 4.8 bis 中宣布。
\begin{SGAresult}{推论}{2.5}
设 $S$ 为\Term{prescheme}{预概形}，$G$ 为在 $S$ 上\Term{smooth}{光滑}且\Term{finite-presentation}{有限表示}、\Term{fiber}{纤维}\Term{connected}{连通}的 $S$-\Term{group-prescheme}{群预概形}，$P$ 为 $G$ 的一个\Term{parabolic-subgroup}{抛物子群}（第 XIV 讲 4.8 bis）。则 $G/P$ 由一个在 $S$ 上\Term{smooth}{光滑}且\Term{projective}{射影}的 $S$-\Term{prescheme}{预概形}表示。
\end{SGAresult}
\begin{proof}
只需证明（同上）$G/P$ 由一个\Term{quasi-projective}{拟射影}、\Term{finite-presentation}{有限表示}的 $S$-\Term{prescheme}{预概形}表示；这由 2.4 得到，因为 $P=\operatorname{Norm}_G(P)$（同上），故在 $S$ 上\Term{flat}{平坦}。
\end{proof}

\SGAsection{3}{中心平坦的群}
\begin{SGAresult}{命题}{3.1}
设 $S$ 为\Term{prescheme}{预概形}，$G$ 为在 $S$ 上\Term{smooth}{光滑}且\Term{finite-presentation}{有限表示}、\Term{fiber}{纤维}\Term{connected}{连通}的 $S$-\Term{group-prescheme}{群预概形}。假设 $G$ 的\Term{center}{中心} $Z$（由第 XI 讲 6.11，它\Term{representable}{可表示}）在 $S$ 上\Term{flat}{平坦}。对每个整数 $n>0$，以 $G^{(n)}$ 记等于 $G$ 沿单位\Term{section}{截面}的第 $n$ 个\Term{normal-invariant}{法不变量}的\Term{locally-free-module}{局部自由} $\mathscr O_S$-模，以 $\rho_n$ 记 $G$ 在 $\operatorname{GL}_S(G^{(n)})$ 中自然的“\Term{adjoint-representation}{伴随表示}”。则：

 a）商 $G'=G/Z$ 由一个在 $S$ 上\Term{smooth}{光滑}、\Term{finite-presentation}{有限表示}、\Term{quasi-affine}{拟仿射}且\Term{fiber}{纤维}\Term{affine}{仿射}\Term{connected}{连通}的 $S$-\Term{group-prescheme}{群预概形}表示。

 b）表示 $\rho_n$ 通过 $G'$ \Term{resolution}{分解}：
\[
\begin{tikzcd}
 G\arrow[rr,"\rho_n"]\arrow[dr]&&\operatorname{GL}_S(G^{(n)})\\
 &G/Z\arrow[ur,"i_n"']&
\end{tikzcd}
\]
若 $S$ \Term{quasi-compact}{拟紧}，则当 $n$ 充分大时，$i_n$ 是\Term{monomorphism}{单态射}。

 c）$G'$ 的\Term{subtorus}{子环面}\Term{functor}{函子} $\mathscr T_{G'}$ 由一个在 $S$ 上\Term{smooth}{光滑}的 $S$-\Term{prescheme}{预概形}表示；若 $S$ \Term{quasi-compact}{拟紧}，它还是一族在 $S$ 上\Term{affine}{仿射}的\Term{prescheme}{预概形}之和。
\end{SGAresult}
\begin{proof}
a）中的断言在 $S$ 上是\Term{local-scheme}{局部}的，故可假设 $S$ \Term{quasi-compact}{拟紧}。由第 XI 讲 6.11 b），当 $n$ 充分大时，$Z=\operatorname{Ker}\rho_n$，故后者在 $S$ 上\Term{flat}{平坦}。于是由 2.3，$G/Z$ \Term{representable}{可表示}且 $i_n$ 为\Term{monomorphism}{单态射}。由于 $G$ 在 $S$ 上\Term{smooth}{光滑}且\Term{finite-presentation}{有限表示}，
\SourcePage{336}
$G'$ 也如此（第 VI\textsubscript{B} 讲 §9）。群 $\operatorname{GL}_S(G^{(n)})$ 在 $S$ 上\Term{affine}{仿射}，而\Term{finite-presentation}{有限表示}的\Term{monomorphism}{单态射}是\Term{quasi-affine}{拟仿射}的（EGA IV 8.11.2），故 $G'$ 在 $S$ 上\Term{quasi-affine}{拟仿射}且\Term{fiber}{纤维}\Term{affine}{仿射}。因 $G'$ 在 $S$ 上\Term{smooth}{光滑}，\Term{functor}{函子} $\mathscr T_{G'}$ \Term{formally-smooth}{形式光滑}（第 XI 讲 2.1 bis）。结合第 XI 讲 4.6 和 4.3，3.1 c）中的断言将由下一个\Term{lemma}{引理}得到。
\end{proof}
\begin{SGAresult}{\Term{lemma}{引理}}{3.2}
设 $S$ 为\Term{prescheme}{预概形}，$G,H$ 为两个在 $S$ 上\Term{finite-presentation}{有限表示}的 $S$-\Term{group-prescheme}{群预概形}，$u:G\to H$ 为 $S$-群\Term{monomorphism}{单态射}，$\mathscr T_G$、$\mathscr T_H$ 分别为 $G$、$H$ 的\Term{subtorus}{子环面}\Term{functor}{函子}（参见第 XV 讲 §8）。则映射 $T\mapsto u(T)$ 定义一个\Term{monomorphism}{单态射}（第 XV 讲 8.3 c））
\[\widetilde u:\mathscr T_G\longrightarrow\mathscr T_H,\]
它由一个\Term{finite-presentation}{有限表示}的\Term{closed-immersion}{闭浸入}表示。
\end{SGAresult}
\begin{proof}
由通常方法，归约到以下问题：设 $T$ 为 $H$ 的一个\Term{subtorus}{子环面}，$T'=G\times_T H$ 为它在 $G$ 中的\Term{inverse-image}{逆像}，$u_T:T'\to T$ 为由 $u$ 得到的 $S$-群\Term{monomorphism}{单态射}。必须证明 $S$-\Term{functor}{函子} $\prod_{T/S}(T'/T)$ 由 $S$ 的一个\Term{finite-presentation}{有限表示}的\Term{closed-subprescheme}{闭子预概形}表示。但 $T$ 是\Term{torus}{环面}，故在 $S$ 上\Term{smooth}{光滑}、\Term{fiber}{纤维}\Term{connected}{连通}，只需应用第 XI 讲 6.10。
\end{proof}
\begin{SGAresult}{\Term{theorem}{定理}}{3.3}
设 $S$ 为\Term{prescheme}{预概形}，$G$ 为在 $S$ 上\Term{smooth}{光滑}、\Term{finite-presentation}{有限表示}且\Term{fiber}{纤维}\Term{connected}{连通}的 $S$-\Term{group-prescheme}{群预概形}，$Z$ 为 $G$ 的\Term{center}{中心}。假设 $Z$ 在 $S$ 上\Term{flat}{平坦}，并且 $G$ 的\Term{unipotent-rank}{幺幂秩}（第 XV 讲 6.1 ter）等于 $Z$ 的\Term{unipotent-rank}{幺幂秩}。则：

 a）群 $G'=G/Z$ \Term{representable}{可表示}；若 $S$ \Term{quasi-compact}{拟紧}，则当 $n$ 充分大时，\Term{canonical}{典范}\Term{morphism}{态射} $i_n:G'\to\operatorname{GL}_S(G^{(n)})$（参见 3.1 b））为\Term{immersion}{浸入}。

 b）群 $G'$ 在 $S$ 上\Term{quasi-affine}{拟仿射}、\Term{fiber}{纤维}\Term{affine}{仿射}，$G'$ 的\Term{center}{中心}是\Term{unit-group}{单位群}，且 $G'$ 的\Term{unipotent-rank}{幺幂秩}为零。

 c）群 $G'$ 对\Term{etale-topology}{平展拓扑}\Term{local-scheme}{局部}具有\Term{maximal-torus}{极大环面}，它们也都是 $G'$ 的\Term{cartan-subgroup}{Cartan 子群}。$G'$ 的\Term{maximal-torus}{极大环面}\Term{functor}{函子} $\mathscr{TM}_{G'}$（第 XV 讲 §8）由一个在 $S$ 上\Term{smooth}{光滑}且\Term{affine}{仿射}的 $S$-\Term{prescheme}{预概形}表示。

 d）群 $G$ 对\Term{etale-topology}{平展拓扑}\Term{local-scheme}{局部}具有\Term{cartan-subgroup}{Cartan 子群}；$G$ 的\Term{cartan-subgroup}{Cartan 子群}\Term{functor}{函子} $\mathscr C_G$ 由一个在 $S$ 上\Term{smooth}{光滑}且\Term{affine}{仿射}的 $S$-\Term{prescheme}{预概形}表示。
\end{SGAresult}
\begin{proof}
由 3.1，群 $G'$ 由一个在 $S$ 上\Term{smooth}{光滑}、\Term{quasi-affine}{拟仿射}且\Term{fiber}{纤维}\Term{affine}{仿射}的 $S$-\Term{prescheme}{预概形}表示。由 $G$ 和 $G'$ 的各\Term{fiber}{纤维}的\Term{cartan-subgroup}{Cartan 子群}之间的对应（第 XII 讲 6.6 e）），关于 $Z$ 的\Term{unipotent-rank}{幺幂秩}的假设蕴涵 $G'$ 的\Term{unipotent-rank}{幺幂秩}为零。再用第 XV 讲 8.18，可知 $G'$ 对\Term{etale-topology}{平展拓扑}\Term{local-scheme}{局部}具有\Term{maximal-torus}{极大环面}。由 3.1 b）和 1.4 a），当 $n$ 充分大时，$i_n$ 是\Term{immersion}{浸入}。a）得证。

现证明 c）。由于 $G'$ 的\Term{unipotent-rank}{幺幂秩}为零，显然 $G'$ 的每个\Term{maximal-torus}{极大环面}也是 $G'$ 的\Term{cartan-subgroup}{Cartan 子群}。$G'$ 的\Term{maximal-torus}{极大环面}\Term{functor}{函子}由 $G'$ 的\Term{subtorus}{子环面}\Term{functor}{函子}的一个既开又闭的\Term{subprescheme}{子预概形}表示，且它在 $S$ 上\Term{finite-type}{有限型}（第 XV 讲 8.15）；
\SourcePage{337}
于是由 3.1 c），此\Term{functor}{函子}由一个在 $S$ 上\Term{smooth}{光滑}且\Term{affine}{仿射}的 $S$-\Term{prescheme}{预概形}表示。

由于 $G'$ 对\Term{etale-topology}{平展拓扑}\Term{local-scheme}{局部}具有\Term{cartan-subgroup}{Cartan 子群}，$G$ 也如此；而 $G$ 的\Term{cartan-subgroup}{Cartan 子群}\Term{functor}{函子}\Term{canonical-isomorphism}{典范同构}于 $G'$ 的相应\Term{functor}{函子}（第 XV 讲 7.3 iv）），故 c）$\Rightarrow$ d）。

还需证明 b），更确切地说，需证明 $G'$ 的\Term{center}{中心} $Z'$ 是\Term{unit-group}{单位群}。由于 $Z'$ \Term{representable}{可表示}（第 XI 讲 6.11）且在 $S$ 上\Term{finite-presentation}{有限表示}，只需证明 $Z'$ 的\Term{fiber}{纤维}均为\Term{unit-group}{单位群}（第 VI\textsubscript{B} 讲 2.9），因而归约到 $S$ 是\Term{algebraically-closed}{代数闭}域之\Term{spectrum}{谱}的情形。\Term{center}{中心} $Z'$ 显然包含在 $G'$ 的每个\Term{cartan-subgroup}{Cartan 子群}中，从而由 c）包含在 $G'$ 的每个\Term{maximal-torus}{极大环面}中，故为\Term{multiplicative-type}{乘法型}。另一方面，我们将在第 XVII 讲中看到，若 $G$ 是\Term{center}{中心}为 $Z$ 的\Term{connected}{连通}\Term{algebraic-group}{代数群}，则 $G/Z$ 的\Term{center}{中心} $Z'$ 为\Term{unipotent}{幺幂}群。在这里，$Z'$ 既为\Term{multiplicative-type}{乘法型}又为\Term{unipotent}{幺幂}群，故为\Term{unit-group}{单位群}（参见第 XVII 讲）。
\end{proof}
\textbf{\Term{center}{中心}\Term{flat}{平坦}的群的例子。}
\begin{SGAresult}{命题}{3.4}
设 $S$ 为\Term{prescheme}{预概形}，$G$ 为在 $S$ 上\Term{finite-presentation}{有限表示}、\Term{smooth}{光滑}且\Term{fiber}{纤维}\Term{connected}{连通}的 $S$-\Term{group-prescheme}{群预概形}，$Z$ 为 $G$ 的\Term{center}{中心}。在下列两种情形中，$Z$ 在 $S$ 上\Term{flat}{平坦}：

 a）$G$ 的\Term{fiber}{纤维}的\Term{unipotent-rank}{幺幂秩} $\rho_u$（第 XV 讲 6.1 ter）为零。

 b）i）$S$ \Term{reduced}{既约}；ii）$Z$ 的\Term{fiber}{纤维}的\Term{dimension}{维数}是 $S$ 上的\Term{locally-constant}{局部常值}函数；iii）$G$ 的\Term{unipotent-rank}{幺幂秩}等于 $Z$ 的\Term{unipotent-rank}{幺幂秩}。
\end{SGAresult}
\textbf{3.4 a）的证明。}我们将同时证明以下命题。
\begin{SGAresult}{命题}{3.5}
在 3.4 a）的假设下，再假设 $G$ 具有一个\Term{maximal-torus}{极大环面} $T$，并令 $C=\operatorname{Centr}_G(T)$ 为 $G$ 的与 $T$ 相关联的\Term{cartan-subgroup}{Cartan 子群}（第 XII 讲 7.1）。则：

（i）$Z\cap T$ 为 $G$ 的一个\Term{multiplicative-type}{乘法型}\Term{subgroup}{子群}。

（ii）$C$ \Term{commutative-group}{交换}，且等于 $T\cdot Z$。

（iii）若商 $Z/(Z\cap T)$ 和 $C/T$ \Term{representable}{可表示}，则它们由\Term{abelian-prescheme}{阿贝尔预概形}（即在 $S$ 上\Term{smooth}{光滑}、\Term{fiber}{纤维}为\Term{abelian-variety}{阿贝尔簇}的 $S$-\Term{group-prescheme}{群预概形}）表示，且\Term{canonical}{典范}\Term{monomorphism}{单态射} $Z/(Z\cap T)\to C/T$ 是\Term{isomorphism}{同构}。
\end{SGAresult}
利用第 VI\textsubscript{B} 讲 §10 及第 XV 讲 6.2、6.3、6.3 bis 中证明的一般\Term{limit}{极限}性质，照例归约到 $S$ \Term{noetherian}{诺特}的情形（注意 3.4 和 3.5 中的各项断言在 $S$ 上是\Term{local-scheme}{局部}的）。我们已在 3.1 中指出，对 $G$ 的假设蕴涵 $Z$ \Term{representable}{可表示}。于是，由 EGA 0\textsubscript{III} 10.2.6，为证明 $Z$ 在 $S$ 上\Term{flat}{平坦}，可归约到 $S$ 为\Term{local-artinian}{局部阿廷}的情形。这时，由于 $G$ \Term{smooth}{光滑}，$G$ 对\Term{etale-topology}{平展拓扑}\Term{local-scheme}{局部}具有\Term{maximal-torus}{极大环面}（第 XV 讲 8.17）。对 $S$ 作\Term{finite}{有限}\Term{flat}{平坦}\Term{group-extension}{扩张}后（为证明 $Z$ \Term{flat}{平坦}，允许这样做），便可假设 $G$ 具有一个\Term{maximal-torus}{极大环面} $T$。
\SourcePage{338}
同样说明，为证明 3.5，也可归约到 $S$ \Term{artinian}{阿廷}的情形。

 i）由于 $Z$ 在 $G$ 中\Term{closed}{闭}（第 XI 讲 6.11），$Z\cap T$ 是 $T$ 的\Term{closed}{闭}\Term{group-subscheme}{群子概形}。由第 X 讲 4.8 b），$Z\cap T$ 为\Term{multiplicative-type}{乘法型}，当且仅当它在 $S$ 上\Term{flat}{平坦}。同上，只需在 $S$ \Term{artinian}{阿廷}时证明 $Z\cap T$ \Term{flat}{平坦}。

 ii）由于 $G$ 的\Term{unipotent-rank}{幺幂秩}为零，$G$ 的每个\Term{cartan-subgroup}{Cartan 子群} $C$ 都满足\Term{lemma}{引理} 1.6 的假设，故为\Term{commutative-group}{交换群}。$C=T\cdot Z$ 将由 iii）得到。

 iii）若 $C/T$ \Term{representable}{可表示}，则 $C/T$ 在 $S$ 上\Term{smooth}{光滑}（第 VI\textsubscript{B} 讲 §9），且其\Term{fiber}{纤维}为\Term{abelian-variety}{阿贝尔簇}（1.7），故 $C/T$ 为\Term{abelian-prescheme}{阿贝尔预概形}。为证明\Term{canonical}{典范}\Term{monomorphism}{单态射}
\[Z/(Z\cap T)\xrightarrow{i}C/T\]
为\Term{isomorphism}{同构}，只需在 $S$ \Term{local-artinian}{局部阿廷}时验证。事实上，依次可\Term{pushout}{推出} $Z/(Z\cap T)$ 在 $S$ 上\Term{flat}{平坦}（EGA 0\textsubscript{III} 10.2.6），$i$ \Term{flat}{平坦}（EGA IV 11.3.10），$i$ 为\Term{open-immersion}{开浸入}（EGA IV 17.9.1），最后 $i$ 为\Term{isomorphism}{同构}（因为 $C/T$ 的\Term{fiber}{纤维}\Term{connected}{连通}）。

以下假设 $S$ \Term{local-artinian}{局部阿廷}，\Term{closed-point}{闭点}为 $s$，且 $G$ 具有一个\Term{maximal-torus}{极大环面} $T$。令 $C=\operatorname{Centr}_G(T)$。群 $C$ 是 $G$ 的\Term{cartan-subgroup}{Cartan 子群}（第 XII 讲 7.1），且包含 $Z$。

\Term{algebraic-group}{代数群} $M_s=Z_s\cap T_s$ 是 $G_s$ 的一个\Term{central-subgroup}{中心}\Term{multiplicative-type}{乘法型}\Term{subgroup}{子群}。因 $T$ 在 $S$ 上\Term{smooth}{光滑}，$M_s$ \Term{lifting}{提升}为 $T$ 的一个\Term{multiplicative-type}{乘法型}\Term{group-subscheme}{群子概形} $M$（第 IX 讲 3.6 bis），且 $M$ 包含在 $G$ 的\Term{center}{中心}中（第 IX 讲 3.9），从而包含在 $Z\cap T$ 中。由于 $S$ \Term{artinian}{阿廷}，$M$ 和 $T$ 在 $S$ 上\Term{flat}{平坦}，\Term{quotient-group}{商群} $(Z\cap T)/M$、$Z'=Z/M$ 和 $A=C/T$ 都\Term{representable}{可表示}（第 VI\textsubscript{A} 讲 §4 和 §5）。由构造，$((Z\cap T)/M)_s$ 为\Term{unit-group}{单位群}，故 $Z\cap T=M$（第 VI\textsubscript{B} 讲 2.9），特别地，它在 $S$ 上\Term{flat}{平坦}，从而证明了 3.5 i）。

取商后，有一个\Term{canonical}{典范}\Term{monomorphism}{单态射} $i:Z'\to A$，因而是\Term{closed-immersion}{闭浸入}（第 VI\textsubscript{B} 讲 1.4.2）。已经指出，$A$ 是\Term{abelian-scheme}{阿贝尔概形}。还需证明 $i$ 为\Term{isomorphism}{同构}。事实上，这既证明 3.5 iii），也蕴涵 $Z'$ 在 $S$ 上\Term{flat}{平坦}，从而 $Z$ 在 $S$ 上\Term{flat}{平坦}（因为它是 $Z'$ 通过\Term{flat}{平坦}群 $M$ 的\Term{extension}{扩张}，第 VI\textsubscript{B} 讲 §9）。由于 $Z_s$ 和 $G_s$ 的\Term{abelian-rank}{阿贝尔秩}相同（第 XII 讲 6.1），$i_s$ 是\Term{epimorphism}{满态射}，故为\Term{isomorphism}{同构}。取一个在 $S$ 上\Term{invertible}{可逆}的整数 $q>0$。由 1.6 iii）的\Term{dense}{稠密}性\Term{theorem}{定理}，为证明 $i(Z')=A$，只需证明：对每个为 $q$ 的幂的整数 $n$，$i(S')$ 包含 ${}_nA$。记 $M_s^0$ 为 $M_s$ 的\Term{identity-component}{单位分支}。由\Term{duality}{对偶性}立即可见，$M_s^0$ 中的 $q$ 次幂映射为\Term{epimorphism}{满态射}。由此立即\Term{pushout}{推出}：若 $m_0$ 为 $\operatorname{card}(M_s/M_s^0)$ 的素因数\Term{resolution}{分解}中 $q$ 的指数，则 ${}_{nm_0}Z_s$ 在 $A_s\simeq Z_s/M_s$ 中的像包含 ${}_nA_s$。因此存在 $Z_s$ 的一个\Term{multiplicative-type}{乘法型}\Term{subgroup}{子群} $M_s(n)$，其在 $A_s$ 中的像为 ${}_nA_s$。与上文一样，$M_s(n)$ \Term{liftable}{可提升}为 $G$ 的一个\Term{central-subgroup}{中心}\Term{multiplicative-type}{乘法型}\Term{subgroup}{子群} $M(n)$。$M(n)$ 在 $A$ 中的像为\Term{multiplicative-type}{乘法型}\Term{subgroup}{子群}（第 IX 讲
\SourcePage{339}
6.8），故必等于 ${}_nA$，因为在\Term{reduced}{既约}\Term{fiber}{纤维}上如此（第 IX 讲 5.1 bis）。这就完成了 3.4 a）和 3.5 的证明。

\textbf{3.4 b）的证明。}待证断言在 $S$ 上是\Term{local-scheme}{局部}的，故可假设 $S$ \Term{affine}{仿射}，其环为 $A$。将 $A$ 看作其\Term{finite-type}{有限型} $\mathbb Z$-\Term{subalgebra}{子代数}的\Term{inductive-limit}{归纳极限}，同上归约到 $S$ \Term{noetherian}{诺特}且\Term{reduced}{既约}的情形。

为证明 $Z$ \Term{flat}{平坦}，可使用\Term{flatness-valuative-criterion}{平坦性赋值判据}（EGA IV 11.8.1），从而归约到 $S$ 是一个\Term{residue-field}{剩余域}\Term{algebraically-closed}{代数闭}的\Term{complete}{完备}\Term{discrete-valuation-ring}{离散赋值环}之\Term{spectrum}{谱}的情形，记其\Term{generic-point}{泛点}为 $t$、\Term{closed-point}{闭点}为 $s$。令 $Z'$ 为 $Z_t$ 在 $Z$ 中的\Term{schematic-closure}{概形闭包}。必须证明 $Z'=Z$，甚至只需证明 $Z_s=Z'_s$（第 VI\textsubscript{B} 讲 2.6）。由第 VI\textsubscript{B} 讲 §4，有
\[\dim Z_t=\dim Z'_t=\dim Z'_s\le\dim Z_s.\]
但由假设，$\dim Z_t=\dim Z_s$，故 $\dim Z'_s=\dim Z_s$，从而 $Z'_s$ 包含 $(Z_s)^0_{\mathrm{red}}$。因此 $G'_s=G_s/Z'_s$ \Term{affine}{仿射}（第 XII 讲 6.1）。由 $G$ 与 $G'$ 的\Term{cartan-subgroup}{Cartan 子群}之间的对应（第 XII 讲 6.6 e）），假设 3.4 iii）蕴涵 $G'_s$ 的\Term{unipotent-rank}{幺幂秩}为零，从而它的\Term{cartan-subgroup}{Cartan 子群}也都是它的\Term{maximal-torus}{极大环面}。$Z_s$ 在 $G'_s$ 中的像 $Z''_s$ 为 $G'_s$ 的\Term{central-subgroup}{中心}\Term{algebraic-subgroup}{代数子群}，故包含在 $G'_s$ 的每个\Term{cartan-subgroup}{Cartan 子群}中，也就是包含在每个\Term{maximal-torus}{极大环面}中；所以 $Z''_s$ 是 $G'_s$ 的\Term{multiplicative-type}{乘法型}\Term{subgroup}{子群}。在这些条件下，我们将在第 XVII 讲 §7 中证明：存在 $Z_s$ 的一个\Term{finite}{有限}\Term{multiplicative-type}{乘法型}\Term{subgroup}{子群} $M_s$，它在 $Z_s/Z'_s$ 中的像为 $Z''_s$（在 3.4 a）的证明中，当 $Z''_s$ \Term{etale}{平展}时，已使用过这一事实）。因 $G$ 在 $S$ 上\Term{smooth}{光滑}且 $S$ 为\Term{complete}{完备}\Term{local-ring}{局部环}之\Term{spectrum}{谱}，$M_s$ \Term{liftable}{可提升}为 $G$ 的一个 $S$-\Term{multiplicative-type}{乘法型}\Term{subgroup}{子群} $M$（第 IX 讲 3.6 bis 和第 XV 讲 1.6 b）），它还是\Term{central-subgroup}{中心}\Term{subgroup}{子群}（第 IX 讲 5.6 a））。故 $M_t$ 包含在 $Z'_t=Z_t$ 中；又因 $M$ \Term{flat}{平坦}，$M$ 包含在 $Z'$ 中。所以 $M_s$ 包含在 $Z'_s$ 中，但这蕴涵 $Z''_s$ 为\Term{unit-group}{单位群}，即 $Z_s=Z'_s$。3.4 b）得证。

\textbf{\Term{center}{中心}不\Term{flat}{平坦}的\Term{smooth}{光滑}且\Term{fiber}{纤维}\Term{connected}{连通}的\Term{group-prescheme}{群预概形}的例子。}

设 $S$ 为\Term{discrete-valuation-ring}{离散赋值环} $A$ 的\Term{spectrum}{谱}，$\pi$ 为 $A$ 的\Term{uniformizer}{一致化元}，$t$ 为 $S$ 的\Term{generic-point}{泛点}，$s$ 为\Term{closed-point}{闭点}。设 $G$ 为在 $S$ 上\Term{smooth}{光滑}且\Term{affine}{仿射}的 $S$-群，其环为 $B=A[T,T^{-1},U]/(F)$，其中 $F=1-T+\pi U$，\Term{composition-law}{合成律}由下式定义：
\[(t,u),(t',u')\longmapsto(tt',\pi uu'+u+u').\]
于是\Term{generic-fiber}{泛纤维} $G_t$ \Term{isomorphism}{同构}于\Term{multiplicative-group}{乘法群} $(\mathbb G_m)_t$，而\Term{closed-fiber}{闭纤维} $G_s$ \Term{isomorphism}{同构}于\Term{additive-group}{加法群} $(\mathbb G_a)_s$。函数 $T$ 在 $\Gamma(G,\mathscr O_G)$ 中\Term{invertible}{可逆}，并定义一个 $S$-\Term{group-morphism}{群态射} $\varphi:G\to(\mathbb G_m)_S$；它在\Term{generic-fiber}{泛纤维}上为\Term{isomorphism}{同构}，在\Term{closed-fiber}{闭纤维}上为零\Term{morphism}{态射}。给定 $\varphi$ 便\Term{constructible}{可构造}\Term{semidirect-product}{半直积}群 $H=G\cdot(\mathbb G_a)_S$（注意 $(\mathbb G_m)_S$ \Term{action}{作用}于 $(\mathbb G_a)_S$）。$H$ 的\Term{center}{中心} $Z$ 在\Term{generic-fiber}{泛纤维}上为\Term{unit-group}{单位群}，在\Term{closed-fiber}{闭纤维}上等于 $H_s$，故在 $S$ 上不\Term{flat}{平坦}。

\SourcePage{340}
\SGAsection{4}{纤维仿射且幺幂秩为零的群}
\begin{SGAresult}{\Term{theorem}{定理}}{4.1}
设 $S$ 为\Term{prescheme}{预概形}，$G$ 为在 $S$ 上\Term{smooth}{光滑}、\Term{finite-presentation}{有限表示}，\Term{fiber}{纤维}\Term{affine}{仿射}、\Term{connected}{连通}且\Term{unipotent-rank}{幺幂秩}（第 XV 讲 6.1 ter）为零的 $S$-\Term{group-prescheme}{群预概形}。则：

 a）$G$ 的\Term{center}{中心} $Z$ 为\Term{multiplicative-type}{乘法型}群（从而是 $G$ 的一个\Term{reductive-center}{约化中心}，第 XII 讲 4.1）。

 b）$G$ 满足 3.3 的条件 a）至 d）。

 c）$G$ 对\Term{etale-topology}{平展拓扑}\Term{local-scheme}{局部}具有\Term{maximal-torus}{极大环面}，且它们也都是 $G$ 的\Term{cartan-subgroup}{Cartan 子群}。$G$ 的\Term{maximal-torus}{极大环面}\Term{functor}{函子} $\mathscr{TM}_G$ 由一个在 $S$ 上\Term{smooth}{光滑}且\Term{affine}{仿射}的 $S$-\Term{prescheme}{预概形}表示。

 d）$G$ 在 $S$ 上\Term{quasi-affine}{拟仿射}。此外，若 $T$ 为 $G$ 的一个\Term{maximal-torus}{极大环面}，$N=\operatorname{Norm}_G T$ 为它在 $G$ 中的\Term{normalizer}{正规化子}，$W=N/T$ 为关于 $T$ 的\Term{weyl-group}{Weyl 群}，则下列条件\Term{equivalence}{等价}：

（i）$G$ 在 $S$ 上\Term{affine}{仿射}。

（ii）$G'=G/Z$ 在 $S$ 上\Term{affine}{仿射}。

（iii）$N$ 在 $S$ 上\Term{affine}{仿射}。

（iv）$W$ 在 $S$ 上\Term{affine}{仿射}。

若 $S$ \Term{locally-noetherian}{局部诺特}且\Term{dimension}{维数} $\le1$，这些条件总是成立。
\end{SGAresult}
\begin{proof}
由于 $G$ 的\Term{fiber}{纤维}\Term{affine}{仿射}（故\Term{abelian-rank}{阿贝尔秩}为零）且\Term{unipotent-rank}{幺幂秩}为零，$G$ 对\Term{etale-topology}{平展拓扑}\Term{local-scheme}{局部}具有\Term{maximal-torus}{极大环面}（第 XV 讲 8.18），而 $G$ 的每个\Term{maximal-torus}{极大环面}显然也是 $G$ 的\Term{cartan-subgroup}{Cartan 子群}。这证明了 c）的第一部分。为证明 $Z$ 为\Term{multiplicative-type}{乘法型}，由上文可假设 $G$ 具有一个\Term{maximal-torus}{极大环面} $T$。由于 $T$ 也是\Term{cartan-subgroup}{Cartan 子群}，$T$ 包含 $Z$（因为由第 XII 讲 7.1，$T=\operatorname{Centr}_G(T)$），而由 3.5 i），$Z=Z\cap T$ 为\Term{multiplicative-type}{乘法型}。这证明 a）。结合 a），断言 b）显然成立。另一方面，\Term{functor}{函子} $\mathscr{TM}_G$ \Term{isomorphism}{同构}于 $G$ 的\Term{cartan-subgroup}{Cartan 子群}\Term{functor}{函子}（第 XII 讲 7.1），故在 $S$ 上\Term{smooth}{光滑}且\Term{affine}{仿射}（3.3 c）），这就完成了 c）的证明。还需证明 d）。

\emph{d）的证明。}由于 $G'$ \Term{quasi-affine}{拟仿射}（3.5 b）），而 $Z$ 为\Term{multiplicative-type}{乘法型}，从而在 $S$ 上\Term{affine}{仿射}，故 $G$ \Term{quasi-affine}{拟仿射}（第 VI\textsubscript{B} 讲 §9）。

 i）$\Leftrightarrow$ ii）。若 $G$ \Term{affine}{仿射}，则由第 IX 讲 2.3，$G'$ \Term{affine}{仿射}。若 $G'$ \Term{affine}{仿射}，则 $G$ 作为\Term{affine}{仿射}群通过\Term{affine}{仿射}群的\Term{extension}{扩张}，也是\Term{affine}{仿射}的（第 VI\textsubscript{B} 讲 §9）。

 i）$\Leftrightarrow$ iii）。若 $G$ \Term{affine}{仿射}，则 $N$ 在 $G$ 中\Term{closed}{闭}（第 XI 讲 6.11 a）），故\Term{affine}{仿射}。另一方面，$G/N$ \Term{isomorphism}{同构}于\Term{functor}{函子} $\mathscr{TM}_G$（\Term{maximal-torus}{极大环面}的\Term{conjugacy}{共轭}性，参见第 XII 讲 7.1 b）），故由 c）在 $S$ 上\Term{affine}{仿射}。因而若 $N$ 在 $S$ 上\Term{affine}{仿射}，$G$ 作为\Term{affine}{仿射}\Term{homogeneous-space}{齐性空间}通过\Term{affine}{仿射}群的“\Term{extension}{扩张}”，也是\Term{affine}{仿射}的（第 VI\textsubscript{B} 讲 §9）。

 iii）$\Leftrightarrow$ iv）。若 $N$ \Term{affine}{仿射}，则 $W=N/T$ \Term{affine}{仿射}（第 IX 讲 2.3）；反向由第 VI\textsubscript{B} 讲 §9 得到。

此外，$N/T$ 由一个在 $S$ 上\Term{etale}{平展}、\Term{separated}{分离}、
\SourcePage{341}
\Term{finite-type}{有限型}的 $S$-\Term{group-prescheme}{群预概形}表示（第 XV 讲 7.1 iv）），故在 $S$ 上\Term{quasi-finite}{拟有限}。将下一个\Term{lemma}{引理}用于 $X=W$、$Y=S$，便得到 d）的最后一项断言。
\end{proof}
\begin{SGAresult}{\Term{lemma}{引理}}{4.2}
设 $S$ 为\Term{dimension}{维数}为 $1$ 的\Term{locally-noetherian}{局部诺特}\Term{prescheme}{预概形}，$X,Y$ 为两个在 $S$ 上\Term{locally-finite-type}{局部有限型}的 $S$-\Term{prescheme}{预概形}，$u:X\to Y$ 为\Term{quasi-finite}{拟有限}且\Term{separated}{分离}的 $S$-\Term{morphism}{态射}。假设对 $S$ 的每个点 $s$，由 $u$ 经\Term{base-change}{基变换} $\operatorname{Spec}\kappa(s)\to S$ 得到的\Term{morphism}{态射} $u_s:X_s\to Y_s$ \Term{finite}{有限}。则 $u$ 是\Term{affine}{仿射}\Term{morphism}{态射}。
\end{SGAresult}
\begin{proof}
待证断言在 $Y$ 上是\Term{local-scheme}{局部}的，故可假设 $Y$（从而 $X$）在 $S$ 上\Term{finite-type}{有限型}。由 EGA IV 8，只需在 $S$ 为\Term{local-ring}{局部环} $A$ 的\Term{spectrum}{谱}时证明 4.2。由 \Term{fpqc-descent}{fpqc 下降}，可假设 $A$ \Term{complete}{完备}，继而可假设 $A$ \Term{reduced}{既约}（EGA II 1.6.4），再假设 $A$ \Term{normal}{正规}（Nagata \Term{theorem}{定理}，EGA 0\textsubscript{IV} 22；以及 Chevalley \Term{theorem}{定理}，EGA II 6.7.1）。若 $A$ 为域，则由假设 $u$ \Term{finite}{有限}，故\Term{affine}{仿射}。否则，$A$ 为\Term{discrete-valuation-ring}{离散赋值环}；记 $s$ 为 $S$ 的\Term{closed-point}{闭点}，$t$ 为\Term{generic-point}{泛点}，$\pi$ 为 $A$ 的\Term{uniformizer}{一致化元}。设 $y$ 为 $Y_s$ 的一点。再次应用 EGA IV 8，可用 $\mathscr O_{Y,y}$ 的\Term{spectrum}{谱} $Y'$ 替换 $Y$，用 $X'=X\times_Y Y'$ 替换 $X$；最后可假设 $\mathscr O_{Y,y}$ \Term{complete}{完备}。因 $u$ \Term{quasi-finite}{拟有限}且\Term{separated}{分离}，由 EGA II 6.2.6，$X'$ 为两个\Term{scheme}{概形} $X_1$ 与 $X_2$ 的和，其中 $X_1$ 在 $Y'$ 上\Term{finite}{有限}（故\Term{affine}{仿射}），而 $u(X_2)$ 不包含 $Y'$ 的\Term{closed-point}{闭点} $y$。我们断言，在这些条件下，$X_2$ 不与\Term{closed-fiber}{闭纤维} $X'_s$ 相交。事实上，由假设 $X'_s\to Y'_s$ \Term{finite}{有限}，故此\Term{morphism}{态射}在闭集 $(X_2)_s$ 上的\Term{restriction}{限制}\Term{finite}{有限}，其在 $(Y')_s$ 中的像为闭集。由于该像不包含\Term{local-scheme}{局部}\Term{scheme}{概形} $(Y')_s$ 的\Term{closed-point}{闭点}，$(X_2)_s$ 为空。由于 $u_t$ \Term{finite}{有限}，\Term{morphism}{态射} $X'_t\to Y'_t$ 在 $(X_2)_t=X_2$ 上的\Term{restriction}{限制}\Term{finite}{有限}。另一方面，$Y'_t=Y'_\pi$，故\Term{open-immersion}{开浸入} $Y'_t\to Y'$ \Term{affine}{仿射}。总之，\Term{composition-morphism}{复合态射} $X_2\to Y'_t\to Y'$ \Term{affine}{仿射}，从而\Term{morphism}{态射} $X'=X_1\amalg X_2\to Y'$ 确为\Term{affine}{仿射}。
\end{proof}
\begin{SGAresult}{推论}{4.3}
设 $S$ 为\Term{prescheme}{预概形}，$G$ 为在 $S$ 上\Term{smooth}{光滑}且\Term{finite-presentation}{有限表示}，\Term{fiber}{纤维}\Term{affine}{仿射}、\Term{solvable}{可解}、\Term{connected}{连通}且\Term{unipotent-rank}{幺幂秩}为零的 $S$-\Term{group-prescheme}{群预概形}。则 $G$ 在 $S$ 上\Term{affine}{仿射}。若再有 $G$ 的\Term{center}{中心}为\Term{unit-group}{单位群}，且 $S$ \Term{quasi-compact}{拟紧}，则当 $n$ 充分大时，\Term{canonical}{典范}\Term{morphism}{态射} $i_n:G\to\operatorname{GL}_S(G^{(n)})$（参见 3.1 b））是\Term{closed-immersion}{闭浸入}。
\end{SGAresult}
\begin{proof}
为证明 $G$ 在 $S$ 上\Term{affine}{仿射}，可假设 $G$ 具有一个\Term{maximal-torus}{极大环面} $T$（4.1 c））。群 $G$ 的\Term{fiber}{纤维}\Term{solvable}{可解}，故关于 $T$ 的\Term{weyl-group}{Weyl 群}为\Term{unit-group}{单位群}（BIBLE，6，\Term{theorem}{定理} 1，推论 3），条件 4.1 d）iv）成立。若 $G$ 的\Term{center}{中心}为\Term{unit-group}{单位群}，则当 $n$ 充分大时，$i_n$ 是\Term{monomorphism}{单态射}（3.1），故为\Term{closed-immersion}{闭浸入}（1.5 b））。
\end{proof}

\SGAsection{5}{对约化群和半单群的应用}
\begin{SGAresult}{定义}{5.1}
一个 $S$-\Term{group-prescheme}{群预概形} $G$ 称为\Term{reductive-group}{约化群}（相应地，\Term{semisimple-group}{半单群}），如果 $G$ 在 $S$ 上\Term{smooth}{光滑}且\Term{affine}{仿射}，其\Term{fiber}{纤维}\Term{reductive}{约化}（相应地，\Term{semisimple}{半单}）。
\end{SGAresult}
\SourcePage{342}
\textbf{5.1.1.}从第 XIX 讲起，将系统地研究\Term{reductive-group}{约化群}。本节需要以下性质，它们将在第 XIX 讲中证明（不使用本讲的论述）。

 a）设 $S$ 为\Term{prescheme}{预概形}，$G$ 为在 $S$ 上\Term{smooth}{光滑}、\Term{affine}{仿射}且\Term{fiber}{纤维}\Term{connected}{连通}的 $S$-\Term{group-prescheme}{群预概形}，$s$ 为 $S$ 的一点，使 $G_s$ \Term{reductive}{约化}（相应地，\Term{semisimple}{半单}）。则存在 $s$ 的一个\Term{neighborhood}{邻域} $U$，使 $G|_U$ \Term{reductive}{约化}（相应地，\Term{semisimple}{半单}）。

 b）若 $G$ 为\Term{reductive-group}{约化群}，则 $G$ 对\Term{etale-topology}{平展拓扑}\Term{local-scheme}{局部}具有\Term{maximal-torus}{极大环面}；若 $T$ 为 $G$ 的一个\Term{maximal-torus}{极大环面}，则\Term{weyl-group}{Weyl 群} $W=\operatorname{Norm}_G(T)/T$ 在 $S$ 上\Term{finite}{有限}。

 c）\Term{reductive-group}{约化群}的\Term{unipotent-rank}{幺幂秩}为零。

承认这些性质后，我们将改进上面的断言 a）。
\begin{SGAresult}{\Term{theorem}{定理}}{5.2}
设 $S$ 为\Term{prescheme}{预概形}，$G$ 为在 $S$ 上\Term{smooth}{光滑}且\Term{finite-presentation}{有限表示}、\Term{fiber}{纤维}\Term{connected}{连通}的 $S$-\Term{group-prescheme}{群预概形}，$s$ 为 $S$ 的一点。则：

（i）若 $G$ 的\Term{fiber}{纤维}\Term{affine}{仿射}且 $G_s$ \Term{reductive}{约化}，则存在 $s$ 的一个\Term{open-neighborhood}{开邻域} $U$，使 $G|_U$ \Term{reductive}{约化}。

（ii）若 $G_s$ \Term{semisimple}{半单}，则存在 $s$ 的一个\Term{open-neighborhood}{开邻域} $U$，使 $G|_U$ \Term{semisimple}{半单}。
\end{SGAresult}
\begin{proof}
利用第 XV 讲 6.2 i）和第 VI\textsubscript{B} 讲 §10，归约到 $S$ \Term{noetherian}{诺特}的情形。

在 ii）中，考虑 $G$ 的\Term{center}{中心} $Z$，它\Term{representable}{可表示}（第 XI 讲 6.11 a））。由于 $G_s$ \Term{semisimple}{半单}，熟知 $Z_s$ \Term{finite}{有限}。因此（第 VI\textsubscript{B} 讲 §4），存在 $s$ 的一个\Term{neighborhood}{邻域} $U$，使 $Z$ 在 $U$ 上\Term{quasi-finite}{拟有限}。于是对 $U$ 的每个点 $t$，$G_t$ \Term{affine}{仿射}（第 XII 讲 6.1）。因此缩小 $S$ 后，在 ii）中也可像在 i）中一样假设 $G$ 的\Term{fiber}{纤维}\Term{affine}{仿射}。

结合 a），为证明 5.2，只需证明 $G$ 在 $S$ 上\Term{affine}{仿射}。由于 $G$ 的\Term{fiber}{纤维}\Term{affine}{仿射}，已知 $G$ 的\Term{subtorus}{子环面}\Term{functor}{函子}由一个\Term{smooth}{光滑} $S$-\Term{prescheme}{预概形}表示（第 XV 讲 8.11 和 8.9）。作\Term{covering}{覆盖} $s$ 的\Term{etale}{平展}\Term{group-extension}{扩张}后（这是允许的），便可假设存在 $G$ 的一个\Term{subtorus}{子环面} $T$，使 $T_s$ 为 $G_s$ 的\Term{maximal-torus}{极大环面}。于是 $C=\operatorname{Centr}_G(T)$ 为 $G$ 的一个\Term{smooth}{光滑}、\Term{fiber}{纤维}\Term{connected}{连通}且包含 $T$ 的\Term{subgroup}{子群}，满足 $C_s=T_s$（因为 $C_s$ 是 $G_s$ 的\Term{cartan-subgroup}{Cartan 子群}，且由 c），$G_s$ 的\Term{unipotent-rank}{幺幂秩}为零）。由此得 $C=T$，故 $T$ 是 $G$ 的\Term{cartan-subgroup}{Cartan 子群}，也是\Term{maximal-torus}{极大环面}；特别地，$G$ 的\Term{unipotent-rank}{幺幂秩}为零，可以应用 4.1。

由 4.1 d），只需证明关于 $T$ 的\Term{weyl-group}{Weyl 群} $W$ 在 $s$ 的某个\Term{neighborhood}{邻域} $U$ 上\Term{affine}{仿射}。实际上，我们将看到 $W$ 在 $s$ 的某个\Term{neighborhood}{邻域}上甚至\Term{finite}{有限}。由于 $W$ 在 $S$ 上\Term{quasi-finite}{拟有限}（第 XV 讲 7.1 iv）），$W$ 在 $s$ 的某个\Term{neighborhood}{邻域} $U$ 上\Term{finite}{有限}，\Term{equivalence}{等价}于\Term{morphism}{态射} $W\to S$ 在 $s$ 处\Term{proper}{固有}（EGA IV 15.7 和 EGA III 4.4.2）；为证明这一点，可使用\Term{local-properness-valuative-criterion}{局部固有性赋值判据}（EGA IV 15.7），它将我们归约到 $S$ 是\Term{closed-point}{闭点}为 $s$ 的\Term{discrete-valuation-ring}{离散赋值环}之\Term{spectrum}{谱}的情形。此时由 4.1 d）的最后一项断言，$G$ 在 $S$ 上\Term{affine}{仿射}，因而可应用前面回顾的性质 a），得出 $G$ \Term{reductive}{约化}（相应地，\Term{semisimple}{半单}）。现在再利用
\SourcePage{343}
性质 b），得到 $W$ 确实在 $S$ 上\Term{finite}{有限}，这就完成了 5.2 的证明。
\end{proof}

\SGAsection{6}{应用：将环面的某些刚性性质推广到幺幂秩为零的群}
\begin{SGAresult}{命题}{6.1}
设 $S$ 为\Term{prescheme}{预概形}，$G$ 为在 $S$ 上\Term{smooth}{光滑}、\Term{finite-presentation}{有限表示}、\Term{fiber}{纤维}\Term{connected}{连通}\Term{affine}{仿射}且\Term{center}{中心}为\Term{multiplicative-type}{乘法型}的 $S$-\Term{group-prescheme}{群预概形}（例如 $G$ 的\Term{unipotent-rank}{幺幂秩}为零，参见 4.1 a））。则 $G$ 的每个在 $S$ 上\Term{finite-presentation}{有限表示}且\Term{quasi-finite}{拟有限}的\Term{closed}{闭}\Term{normal-subgroup}{正规子群} $K$ 都在 $S$ 上\Term{finite}{有限}。
\end{SGAresult}
\begin{proof}
事实上，对 $S$ 上的每个\Term{geometric-point}{几何点} $s$，$(K_s)_{\mathrm{red}}$ 是 $G_s$ 的\Term{finite}{有限}\Term{etale}{平展}\Term{subgroup}{子群}，它\Term{normal-subgroup}{正规}，故为\Term{central-subgroup}{中心}\Term{subgroup}{子群}（因为 $G_s$ \Term{connected}{连通}）。因此 $K'=Z\cap K$（其中 $Z$ 为 $G$ 的\Term{center}{中心}）与 $K$ 具有相同的底\Term{space}{空间}，只需证明 $K'$ 在 $S$ 上\Term{finite}{有限}。但 $K'$ 是\Term{multiplicative-type}{乘法型}群 $Z$ 的一个\Term{quasi-finite}{拟有限}\Term{closed-subgroup}{闭子群}，故\Term{finite}{有限}，这一点立即可见（在 $S$ 上\Term{local-scheme}{局部}地，对适当的整数 $n$，$K'$ 包含在 ${}_nZ$ 中，而 ${}_nZ$ 在 $S$ 上\Term{finite}{有限}）。
\end{proof}
\begin{SGAresult}{推论}{6.2}
设 $S,G$ 如上，$K$ 为 $G$ 的一个在 $S$ 上\Term{finite-presentation}{有限表示}、在 $G$ 中\Term{normal-subgroup}{正规}且\Term{closed}{闭}的\Term{group-subprescheme}{群子预概形}，$s$ 为 $S$ 的一点。若 $K_s$ \Term{finite}{有限}（相应地，为\Term{unit-group}{单位群}），则存在 $s$ 的一个\Term{open-neighborhood}{开邻域} $U$，使 $K|_U$ 在 $S$ 上\Term{finite}{有限}（相应地，为\Term{unit-group}{单位群}）。
\end{SGAresult}
\begin{proof}
由 $K$ 的\Term{fiber}{纤维}\Term{dimension}{维数}的\Term{upper-semicontinuity}{上半连续性}（第 VI\textsubscript{B} 讲 §4），缩小 $S$ 后，已可假设 $K$ 在 $S$ 上\Term{quasi-finite}{拟有限}，从而\Term{finite}{有限}（6.1）。若再有 $K_s$ 为\Term{unit-group}{单位群}，则由 \Term{nakayama-lemma}{Nakayama 引理}容易得到 $K$ 在 $\operatorname{Spec}\mathscr O_s$ 上为\Term{unit-group}{单位群}，从而在 $s$ 的某个\Term{neighborhood}{邻域}上也为\Term{unit-group}{单位群}。
\end{proof}
\begin{SGAresult}{命题}{6.3}
设 $S$ 为\Term{prescheme}{预概形}，$G$ 为在 $S$ 上\Term{smooth}{光滑}且\Term{finite-presentation}{有限表示}，\Term{fiber}{纤维}\Term{affine}{仿射}、\Term{connected}{连通}且\Term{unipotent-rank}{幺幂秩}为零的 $S$-\Term{group-prescheme}{群预概形}，$H$ 为 $S$-\Term{group-prescheme}{群预概形}，$s$ 为 $S$ 的一点，$u:G\to H$ 为一个 $S$-\Term{group-morphism}{群态射}，使 $\operatorname{Ker}u_s$ 为\Term{central-subgroup}{中心}\Term{subgroup}{子群}。再假设 $H$ 在 $S$ 上\Term{finite-presentation}{有限表示}，或者 $S$ \Term{locally-noetherian}{局部诺特}且 $G$ 在 $S$ 上\Term{locally-finite-type}{局部有限型}。则存在 $s$ 的一个\Term{open-neighborhood}{开邻域} $U$，使得若 $K=\operatorname{Ker}u$，则 $K|_U$ 为\Term{central-subgroup}{中心}\Term{multiplicative-type}{乘法型}\Term{subgroup}{子群}。此外，$G'=(G|_U)/(K|_U)$ \Term{representable}{可表示}，且由 $u$ 取商得到的\Term{morphism}{态射} $u':G'\to H|_U$ 为\Term{immersion}{浸入}。
\end{SGAresult}
\begin{proof}
设 $Z$ 为 $G$ 的\Term{center}{中心}，$K'=K\cap Z$。

 a）$K'$ 在 $s$ 的某个\Term{neighborhood}{邻域} $U$ 上是\Term{multiplicative-type}{乘法型}\Term{subgroup}{子群}。事实上，在 $s$ 的某个\Term{neighborhood}{邻域} $U$ 上作\Term{etale}{平展}\Term{group-extension}{扩张}后（这是允许的），可假设 $G$ 具有一个\Term{maximal-torus}{极大环面} $T$（4.1 c））。于是 $T\cap K$ 是一个\Term{multiplicative-type}{乘法型}群（第 XV 讲 8.3），由假设它在 $s$ 处的\Term{fiber}{纤维}为\Term{central-subgroup}{中心}\Term{subgroup}{子群}，故 $T\cap K$ 在 $s$ 的某个\Term{neighborhood}{邻域}上为\Term{central-subgroup}{中心}\Term{subgroup}{子群}（第 IX 讲 5.6 a）），从而等于 $K'$。
\SourcePage{344}
 b）证明在 $s$ 的某个\Term{neighborhood}{邻域} $U$ 上有 $K'=K$。由 a），已可假设 $K'$ 为\Term{multiplicative-type}{乘法型}，故在 $S$ 上\Term{flat}{平坦}。由于 $K'_s=K_s$，自然\Term{immersion}{浸入} $K'\to K$ 在 $s$ 的某个\Term{neighborhood}{邻域} $U$ 上是开的（第 VI\textsubscript{B} 讲 2.6）。若 $t\in U$，则 $K_t$ 在群 $G'_t=G_t/Z_t$ 中的像为 $G'_t$ 的一个\Term{finite}{有限}\Term{etale}{平展}\Term{normal-subgroup}{正规子群}（因为 $K_t$ 在 $G_t$ 中\Term{normal-subgroup}{正规}），故为\Term{central-subgroup}{中心}\Term{subgroup}{子群}，从而为\Term{unit-group}{单位群}（3.3 b））。即 $K_t$ 包含在 $Z_t$ 中，因而等于 $K'_t$，故在 $U$ 上有 $K=K'$。

 c）$G'$ 的\Term{representability}{可表示性}由 2.3 得到；将 4.1 c）应用于 $G'$ 后，由 1.3 a）可知 $u'$ 为\Term{immersion}{浸入}。
\end{proof}
\begin{SGAresult}{命题}{6.4}
设 $S$ 为\Term{prescheme}{预概形}，$G$ 为在 $S$ 上\Term{smooth}{光滑}、\Term{finite-presentation}{有限表示}且\Term{fiber}{纤维}\Term{connected}{连通}的 $S$-\Term{group-prescheme}{群预概形}，$s$ 为 $S$ 的一点，使 $G_s$ 由其\Term{subtorus}{子环面}生成（第 XII 讲 8.2）（例如 $G_s$ \Term{affine}{仿射}且\Term{unipotent-rank}{幺幂秩}为零），$H$ 为 $S$-\Term{group-prescheme}{群预概形}，$u,v:G\to H$ 为两个 $S$-\Term{homomorphism}{同态}，满足 $u_s=v_s$ 且 $u_s$ 为\Term{central-homomorphism}{中心同态}。再假设 $H$ 在 $S$ 上\Term{finite-presentation}{有限表示}，或者 $S$ \Term{locally-noetherian}{局部诺特}。则存在 $s$ 的一个\Term{neighborhood}{邻域} $U$，使 $u|_U=v|_U$。
\end{SGAresult}
\begin{proof}
照例归约到 $S$ \Term{noetherian}{诺特}的情形（为考察“$G_s$ 由其\Term{subtorus}{子环面}生成”这一条件，使用第 VI\textsubscript{B} 讲 7.4）。由于 $G$ 的\Term{fiber}{纤维}\Term{connected}{连通}，为证明在 $s$ 的某个\Term{neighborhood}{邻域} $U$ 上 $u=v$，只需证明模去\Term{local-ring}{局部环} $\mathscr O_{S,s}$ 的\Term{maximal-ideal}{极大理想}的每个幂后均有 $u=v$，从而归约到 $S$ \Term{local-artinian}{局部阿廷}的情形。

这时，$G$ 的\Term{maximal-subtorus}{极大子环面}\Term{functor}{函子}由一个在 $S$ 上\Term{smooth}{光滑}、\Term{finite-type}{有限型}的 $S$-\Term{scheme}{概形} $X$ 表示（第 XV 讲 8.17）。设 $T$ 为 $G_X$ 的\Term{maximal-subtorus}{极大子环面}，它对于\Term{functor}{函子} $X$ 是\Term{universal}{普遍}的。对 $G_s$ 的假设意即：由 $\kappa(s)$-\Term{morphism}{态射} $f:T_s\to G_s$（即\Term{immersion}{浸入} $T_s\to G_{X_s}$ 与\Term{canonical-projection}{典范投影} $G_{X_s}\to G_s$ 的复合）生成的 $G_s$ 的\Term{algebraic-subgroup}{代数子群}（第 VI\textsubscript{B} 讲 §7）等于整个 $G_s$。而 $X_s$ \Term{geometrically-connected}{几何连通}（因为若 $T$ 是 $G_s$ 的一个\Term{maximal-torus}{极大环面}，$N$ 是它在 $G_s$ 中的\Term{normalizer}{正规化子}，则 $X_s$ \Term{isomorphism}{同构}于 $G_s/N$），且 $T_s$ 在 $f$ 下的像包含 $G_s$ 的单位\Term{section}{截面}；由第 VI\textsubscript{B} 讲 7.4，存在一个整数 $N>0$ 和某个 $S$-\Term{morphism}{态射} $f^N:T^N\to G$（其中 $T^N=T\times_S\cdots\times_S T$，共 $N$ 个因子，而 $f^N$ 仅依赖于 $f$ 和 $G$ 的乘法），使 $(f^N)_s$ \Term{surjectivity}{满射}。

考虑\Term{base-change}{基变换} $X\to S$，以及 $u_X$ 和 $v_X$ 在 $G_X$ 的\Term{subtorus}{子环面} $T$ 上的\Term{restriction}{限制}。由假设，$u_{X_s}=v_{X_s}$，且 $u_{X_s}:T_{X_s}\to H_{X_s}$ 为\Term{central-homomorphism}{中心同态}。于是由第 IX 讲 5.1，$u_X|_T=v_X|_T$。展开 $f^N$ 的定义，并利用 $u,v$ 是\Term{homomorphism}{同态}，立即得到 $uf^N=vf^N$。但 $\kappa(s)$-\Term{morphism}{态射} $f_s^N$ \Term{surjectivity}{满射}，而 $G_s$ \Term{smooth}{光滑}，故\Term{reduced}{既约}，所以 $f_s^N$ \Term{generically-flat}{泛平坦}。由于 $T$ 在 $S$ 上\Term{smooth}{光滑}从而\Term{flat}{平坦}，存在 $T$ 的\Term{nonempty}{非空}\Term{open-subset}{开集} $V$，使 $f^N|_V$ \Term{flat}{平坦}（EGA IV 11.3.10）。$V$ 的像为 $G$ 的一个\Term{open-subset}{开集} $W$（EGA IV 11.3.1），且 $f^N:V\to W$ \Term{faithfully-flat}{忠实平坦}、\Term{finite-presentation}{有限表示}，故为 \Term{fpqc-topology}{fpqc 拓扑}的\Term{covering}{覆盖}。于是等式 $uf^N=vf^N$ 蕴涵 $u|_W=v|_W$。由于 $W$ 在
\SourcePage{345}
$G$ 中\Term{schematically-dense}{概形稠密}（EGA IV 11.10.10），且 $H$ 在 $S$ 上\Term{separated}{分离}（第 VI\textsubscript{A} 讲 0.3），立即得到 $u=v$。
\end{proof}
\SourcePage{346}
% 原 PDF 此页为空白。
